\chapter{Component Families}\label{record-019}

\noindent{\footnotesize\color{BookMuted}\path{docs/design-system/component-families.md}\par}\smallskip

This document is the working semantic map for the shipped Bijou component families.

The goal is to stop thinking in terms of ``every export is its own component.''

\section{How to read this}\label{record-019-how-to-read-this}

\begin{itemize}
\tightlist
\item
  \textbf{Family} means one semantic job.
\item
  \textbf{Variants} may be semantic, interaction-layer, or render-path variants.
\item
  \textbf{Core} means it belongs in \texttt{@\allowbreak{}f\allowbreak{}l\allowbreak{}y\allowbreak{}i\allowbreak{}n\allowbreak{}g\allowbreak{}r\allowbreak{}o\allowbreak{}b\allowbreak{}o\allowbreak{}t\allowbreak{}s\allowbreak{}/\allowbreak{}b\allowbreak{}i\allowbreak{}j\allowbreak{}o\allowbreak{}u}.
\item
  \textbf{TUI} means it belongs in \texttt{@\allowbreak{}f\allowbreak{}l\allowbreak{}y\allowbreak{}i\allowbreak{}n\allowbreak{}g\allowbreak{}r\allowbreak{}o\allowbreak{}b\allowbreak{}o\allowbreak{}t\allowbreak{}s\allowbreak{}/\allowbreak{}b\allowbreak{}i\allowbreak{}j\allowbreak{}o\allowbreak{}u\allowbreak{}-\allowbreak{}t\allowbreak{}u\allowbreak{}i}.
\end{itemize}

\section{Documentation completeness status}\label{record-019-documentation-completeness-status}

This page is a living component-family reference.

The stated design-system standard requires each mature family doc to include:

\begin{itemize}
\tightlist
\item
  what it is
\item
  variations
\item
  when to use
\item
  when not to use
\item
  content guidance
\item
  interaction ownership
\item
  graceful lowering across \texttt{r\allowbreak{}i\allowbreak{}c\allowbreak{}h}, \texttt{s\allowbreak{}t\allowbreak{}a\allowbreak{}t\allowbreak{}i\allowbreak{}c}, \texttt{p\allowbreak{}i\allowbreak{}p\allowbreak{}e}, and \texttt{a\allowbreak{}c\allowbreak{}c\allowbreak{}e\allowbreak{}s\allowbreak{}s\allowbreak{}i\allowbreak{}b\allowbreak{}l\allowbreak{}e}
\item
  related families
\item
  closest Carbon analogue
\end{itemize}

This page now meets the structural documentation standard for shipped families: every family below includes variants, usage guidance, content guidance, ownership, graceful lowering, related families, and a Carbon analogue. The remaining work is about depth and sharper examples, not missing required sections.

\section{Component slice checklist}\label{record-019-component-slice-checklist}

When a component family is changed in code, the matching family docs should be checked in the same slice.

Minimum verification:

\begin{itemize}
\tightlist
\item
  the family is present in this guide
\item
  the listed variants still match reality
\item
  \texttt{U\allowbreak{}s\allowbreak{}e\allowbreak{} \allowbreak{}w\allowbreak{}h\allowbreak{}e\allowbreak{}n} and \texttt{A\allowbreak{}v\allowbreak{}o\allowbreak{}i\allowbreak{}d\allowbreak{} \allowbreak{}w\allowbreak{}h\allowbreak{}e\allowbreak{}n} still match the implementation and examples
\item
  graceful lowering still matches the runtime behavior
\item
  related families still make sense after the code change
\end{itemize}

If those checks fail, the component work is not doctrinally complete yet.

\section{Status and feedback families}\label{record-019-status-and-feedback-families}

\subsection{Inline status}\label{record-019-inline-status}

\begin{itemize}
\tightlist
\item
  Family: \texttt{b\allowbreak{}a\allowbreak{}d\allowbreak{}g\allowbreak{}e\allowbreak{}(\allowbreak{})}
\item
  Variants:

  \begin{itemize}
  \tightlist
  \item
    semantic/status variants
  \end{itemize}
\item
  Use when:

  \begin{itemize}
  \tightlist
  \item
    status is compact and local to another object
  \end{itemize}
\item
  Avoid when:

  \begin{itemize}
  \tightlist
  \item
    message must stand alone or explain itself
  \end{itemize}
\item
  Content guidance:

  \begin{itemize}
  \tightlist
  \item
    keep badge labels terse, stable, and noun-like or state-like
  \item
    avoid embedding full sentences or multi-step guidance in badges
  \end{itemize}
\item
  Ownership:

  \begin{itemize}
  \tightlist
  \item
    core
  \end{itemize}
\item
  Graceful lowering:

  \begin{itemize}
  \tightlist
  \item
    rich/static: keep compact label styling inline with nearby content
  \item
    pipe: lower to plain inline text label without relying on color
  \item
    accessible: speak the status in plain words next to the owning object
  \end{itemize}
\item
  Related families:

  \begin{itemize}
  \tightlist
  \item
    \texttt{a\allowbreak{}l\allowbreak{}e\allowbreak{}r\allowbreak{}t\allowbreak{}(\allowbreak{})}
  \item
    \texttt{n\allowbreak{}o\allowbreak{}t\allowbreak{}e\allowbreak{}(\allowbreak{})}
  \item
    notification system
  \end{itemize}
\item
  Carbon analogue:

  \begin{itemize}
  \tightlist
  \item
    tag
  \end{itemize}
\end{itemize}

\subsection{In-flow status block}\label{record-019-in-flow-status-block}

\begin{itemize}
\tightlist
\item
  Family: \texttt{a\allowbreak{}l\allowbreak{}e\allowbreak{}r\allowbreak{}t\allowbreak{}(\allowbreak{})}
\item
  Variants:

  \begin{itemize}
  \tightlist
  \item
    semantic tones
  \item
    render-path companion: \texttt{a\allowbreak{}l\allowbreak{}e\allowbreak{}r\allowbreak{}t\allowbreak{}S\allowbreak{}u\allowbreak{}r\allowbreak{}f\allowbreak{}a\allowbreak{}c\allowbreak{}e\allowbreak{}(\allowbreak{})}
  \end{itemize}
\item
  Use when:

  \begin{itemize}
  \tightlist
  \item
    status should remain part of the page/document
  \end{itemize}
\item
  Avoid when:

  \begin{itemize}
  \tightlist
  \item
    lifecycle/history/stacking is required
  \end{itemize}
\item
  Content guidance:

  \begin{itemize}
  \tightlist
  \item
    include the state, cause, and next useful action in the body when possible
  \item
    keep the headline short and let supporting text carry detail
  \end{itemize}
\item
  Ownership:

  \begin{itemize}
  \tightlist
  \item
    core
  \end{itemize}
\item
  Graceful lowering:

  \begin{itemize}
  \tightlist
  \item
    rich/static: remain a boxed or otherwise clearly separated in-flow message
  \item
    pipe: lower to emphasized plain text that stays in document order
  \item
    accessible: keep message text and severity explicit, with no dependence on borders or color
  \end{itemize}
\item
  Related families:

  \begin{itemize}
  \tightlist
  \item
    \texttt{b\allowbreak{}a\allowbreak{}d\allowbreak{}g\allowbreak{}e\allowbreak{}(\allowbreak{})}
  \item
    \texttt{n\allowbreak{}o\allowbreak{}t\allowbreak{}e\allowbreak{}(\allowbreak{})}
  \item
    \texttt{t\allowbreak{}o\allowbreak{}a\allowbreak{}s\allowbreak{}t\allowbreak{}(\allowbreak{})}
  \item
    notification system
  \end{itemize}
\item
  Carbon analogue:

  \begin{itemize}
  \tightlist
  \item
    inline notification
  \end{itemize}
\end{itemize}

\subsection{Low-level transient overlay}\label{record-019-low-level-transient-overlay}

\begin{itemize}
\tightlist
\item
  Family: \texttt{t\allowbreak{}o\allowbreak{}a\allowbreak{}s\allowbreak{}t\allowbreak{}(\allowbreak{})}
\item
  Variants:

  \begin{itemize}
  \tightlist
  \item
    status variants
  \item
    anchor/placement
  \end{itemize}
\item
  Use when:

  \begin{itemize}
  \tightlist
  \item
    the app needs a one-off transient overlay and is composing overlays directly
  \item
    placement matters, but lifecycle/history does not
  \end{itemize}
\item
  Avoid when:

  \begin{itemize}
  \tightlist
  \item
    stacking, routing, actions, or recall matter
  \item
    the content should remain in page flow
  \end{itemize}
\item
  Content guidance:

  \begin{itemize}
  \tightlist
  \item
    keep the message short enough to scan without stopping the task
  \item
    avoid multiline operational detail that really belongs in notifications or a drawer
  \end{itemize}
\item
  Ownership:

  \begin{itemize}
  \tightlist
  \item
    TUI
  \end{itemize}
\item
  Graceful lowering:

  \begin{itemize}
  \tightlist
  \item
    rich: anchored transient overlay with placement and timing
  \item
    static: lower to visible in-flow status or last-frame summary rather than hidden motion
  \item
    pipe: lower to one plain event line or app-owned log entry
  \item
    accessible: prefer explicit announcement text over spatial anchoring
  \end{itemize}
\item
  Related families:

  \begin{itemize}
  \tightlist
  \item
    \texttt{a\allowbreak{}l\allowbreak{}e\allowbreak{}r\allowbreak{}t\allowbreak{}(\allowbreak{})}
  \item
    \texttt{m\allowbreak{}o\allowbreak{}d\allowbreak{}a\allowbreak{}l\allowbreak{}(\allowbreak{})}
  \item
    notification system
  \end{itemize}
\item
  Carbon analogue:

  \begin{itemize}
  \tightlist
  \item
    toast notification primitive
  \end{itemize}
\end{itemize}

\subsection{Transient app notifications}\label{record-019-transient-app-notifications}

\begin{itemize}
\tightlist
\item
  Family: notification system
\item
  Variants:

  \begin{itemize}
  \tightlist
  \item
    \texttt{T\allowbreak{}O\allowbreak{}A\allowbreak{}S\allowbreak{}T}
  \item
    \texttt{I\allowbreak{}N\allowbreak{}L\allowbreak{}I\allowbreak{}N\allowbreak{}E}
  \item
    \texttt{A\allowbreak{}C\allowbreak{}T\allowbreak{}I\allowbreak{}O\allowbreak{}N\allowbreak{}A\allowbreak{}B\allowbreak{}L\allowbreak{}E}
  \item
    tones
  \item
    placements
  \item
    history/archive view
  \end{itemize}
\item
  Use when:

  \begin{itemize}
  \tightlist
  \item
    the app owns transient messaging and may need stacking, routing, actions, or recall
  \end{itemize}
\item
  Avoid when:

  \begin{itemize}
  \tightlist
  \item
    the content is primary page content
  \item
    one local overlay is enough and app-wide lifecycle is unnecessary
  \end{itemize}
\item
  Content guidance:

  \begin{itemize}
  \tightlist
  \item
    titles should summarize the event, while body text should explain consequence or next step
  \item
    actionable notifications should expose one clear primary action instead of a menu of choices
  \end{itemize}
\item
  Ownership:

  \begin{itemize}
  \tightlist
  \item
    TUI
  \end{itemize}
\item
  Graceful lowering:

  \begin{itemize}
  \tightlist
  \item
    rich: stacked overlays with placement, action buttons, animation, and archive/history
  \item
    static: visible notices plus accessible history review surface
  \item
    pipe: lower to ordered event lines and routed warnings/errors
  \item
    accessible: simplify to explicit chronological reviewable notices with actions described in text
  \end{itemize}
\item
  Related families:

  \begin{itemize}
  \tightlist
  \item
    \texttt{t\allowbreak{}o\allowbreak{}a\allowbreak{}s\allowbreak{}t\allowbreak{}(\allowbreak{})}
  \item
    \texttt{a\allowbreak{}l\allowbreak{}e\allowbreak{}r\allowbreak{}t\allowbreak{}(\allowbreak{})}
  \item
    \texttt{l\allowbreak{}o\allowbreak{}g\allowbreak{}(\allowbreak{})}
  \item
    \texttt{m\allowbreak{}o\allowbreak{}d\allowbreak{}a\allowbreak{}l\allowbreak{}(\allowbreak{})}
  \end{itemize}
\item
  Carbon analogue:

  \begin{itemize}
  \tightlist
  \item
    toast notification, actionable notification, notification panel/center
  \end{itemize}
\end{itemize}

\subsection{Activity stream}\label{record-019-activity-stream}

\begin{itemize}
\tightlist
\item
  Family: \texttt{l\allowbreak{}o\allowbreak{}g\allowbreak{}(\allowbreak{})}
\item
  Variants:

  \begin{itemize}
  \tightlist
  \item
    log levels
  \end{itemize}
\item
  Use when:

  \begin{itemize}
  \tightlist
  \item
    order and accumulation matter
  \end{itemize}
\item
  Avoid when:

  \begin{itemize}
  \tightlist
  \item
    the message should interrupt instead of accumulate
  \end{itemize}
\item
  Content guidance:

  \begin{itemize}
  \tightlist
  \item
    log lines should remain chronological, specific, and attributable
  \item
    avoid decorative phrasing that makes scanning and filtering harder
  \end{itemize}
\item
  Ownership:

  \begin{itemize}
  \tightlist
  \item
    core
  \end{itemize}
\item
  Graceful lowering:

  \begin{itemize}
  \tightlist
  \item
    rich/static: retain ordered styled lines with level cues
  \item
    pipe: already natural plain sequential output
  \item
    accessible: keep timestamps/levels explicit and avoid decorative formatting
  \end{itemize}
\item
  Related families:

  \begin{itemize}
  \tightlist
  \item
    notification system
  \item
    \texttt{a\allowbreak{}l\allowbreak{}e\allowbreak{}r\allowbreak{}t\allowbreak{}(\allowbreak{})}
  \item
    \texttt{b\allowbreak{}a\allowbreak{}d\allowbreak{}g\allowbreak{}e\allowbreak{}(\allowbreak{})}
  \end{itemize}
\item
  Carbon analogue:

  \begin{itemize}
  \tightlist
  \item
    activity log / event stream pattern
  \end{itemize}
\end{itemize}

\subsection{Loading placeholders}\label{record-019-loading-placeholders}

\begin{itemize}
\tightlist
\item
  Family: \texttt{s\allowbreak{}k\allowbreak{}e\allowbreak{}l\allowbreak{}e\allowbreak{}t\allowbreak{}o\allowbreak{}n\allowbreak{}(\allowbreak{})}
\item
  Variants:

  \begin{itemize}
  \tightlist
  \item
    single-line, multiline, and region-shaped placeholders
  \end{itemize}
\item
  Use when:

  \begin{itemize}
  \tightlist
  \item
    the expected content shape is known but the content is not available yet
  \item
    a short-lived loading interval would otherwise cause distracting layout pop
  \end{itemize}
\item
  Avoid when:

  \begin{itemize}
  \tightlist
  \item
    honest partial content can already be shown
  \item
    the delay is long enough that a clearer loading state, progress indicator, or retry path is needed
  \item
    the placeholder would outlive the user's trust and start to look like fake content
  \end{itemize}
\item
  Content guidance:

  \begin{itemize}
  \tightlist
  \item
    skeletons should roughly match the final information density instead of becoming ornamental gray bars everywhere
  \item
    use them for short-lived uncertainty, then yield to partial or explicit loading states as soon as real content exists
  \item
    keep labels or surrounding context explicit so the user knows what is loading
  \end{itemize}
\item
  Ownership:

  \begin{itemize}
  \tightlist
  \item
    core
  \end{itemize}
\item
  Graceful lowering:

  \begin{itemize}
  \tightlist
  \item
    rich/static: preserve the placeholder shape while loading is genuinely transient
  \item
    pipe: lower to explicit loading text or field labels instead of decorative placeholder bars
  \item
    accessible: announce loading state and the affected region explicitly rather than relying on visual shimmer or shape
  \end{itemize}
\item
  Related families:

  \begin{itemize}
  \tightlist
  \item
    \texttt{s\allowbreak{}p\allowbreak{}i\allowbreak{}n\allowbreak{}n\allowbreak{}e\allowbreak{}r\allowbreak{}(\allowbreak{})}
  \item
    \texttt{p\allowbreak{}r\allowbreak{}o\allowbreak{}g\allowbreak{}r\allowbreak{}e\allowbreak{}s\allowbreak{}s\allowbreak{}B\allowbreak{}a\allowbreak{}r\allowbreak{}(\allowbreak{})}
  \item
    \texttt{n\allowbreak{}o\allowbreak{}t\allowbreak{}e\allowbreak{}(\allowbreak{})}
  \end{itemize}
\item
  Carbon analogue:

  \begin{itemize}
  \tightlist
  \item
    skeleton state
  \end{itemize}
\end{itemize}

\subsection{Inline shortcut cues}\label{record-019-inline-shortcut-cues}

\begin{itemize}
\tightlist
\item
  Family: \texttt{k\allowbreak{}b\allowbreak{}d\allowbreak{}(\allowbreak{})}
\item
  Variants:

  \begin{itemize}
  \tightlist
  \item
    single key, chord, and adjacent grouped key cues
  \end{itemize}
\item
  Use when:

  \begin{itemize}
  \tightlist
  \item
    a local action or instruction needs an inline shortcut hint
  \item
    the shortcut belongs next to the thing it affects
  \end{itemize}
\item
  Avoid when:

  \begin{itemize}
  \tightlist
  \item
    the user needs a broader command reference or grouped shortcut map
  \item
    the chip would become the primary content instead of a supporting cue
  \end{itemize}
\item
  Content guidance:

  \begin{itemize}
  \tightlist
  \item
    show only the keys the user actually needs at that moment
  \item
    keep shortcut chips close to the related action label instead of collecting them into decorative clutter
  \item
    use shell help views for broader keybinding discovery; keep \texttt{k\allowbreak{}b\allowbreak{}d\allowbreak{}(\allowbreak{})} for inline, local hints
  \end{itemize}
\item
  Ownership:

  \begin{itemize}
  \tightlist
  \item
    core
  \end{itemize}
\item
  Graceful lowering:

  \begin{itemize}
  \tightlist
  \item
    rich/static: preserve compact key-chip treatment inline with related content
  \item
    pipe: lower to explicit key names in plain text without decorative framing being the only cue
  \item
    accessible: speak the shortcut and related action in one readable phrase
  \end{itemize}
\item
  Related families:

  \begin{itemize}
  \tightlist
  \item
    \texttt{h\allowbreak{}e\allowbreak{}l\allowbreak{}p\allowbreak{}V\allowbreak{}i\allowbreak{}e\allowbreak{}w\allowbreak{}(\allowbreak{})}
  \item
    \texttt{c\allowbreak{}o\allowbreak{}m\allowbreak{}m\allowbreak{}a\allowbreak{}n\allowbreak{}d\allowbreak{}P\allowbreak{}a\allowbreak{}l\allowbreak{}e\allowbreak{}t\allowbreak{}t\allowbreak{}e\allowbreak{}(\allowbreak{})}
  \item
    \texttt{n\allowbreak{}o\allowbreak{}t\allowbreak{}e\allowbreak{}(\allowbreak{})}
  \end{itemize}
\item
  Carbon analogue:

  \begin{itemize}
  \tightlist
  \item
    keyboard key / shortcut hint
  \end{itemize}
\end{itemize}

\subsection{Progress indicators}\label{record-019-progress-indicators}

\begin{itemize}
\tightlist
\item
  Family:

  \begin{itemize}
  \tightlist
  \item
    \texttt{p\allowbreak{}r\allowbreak{}o\allowbreak{}g\allowbreak{}r\allowbreak{}e\allowbreak{}s\allowbreak{}s\allowbreak{}B\allowbreak{}a\allowbreak{}r\allowbreak{}(\allowbreak{})}
  \item
    \texttt{c\allowbreak{}r\allowbreak{}e\allowbreak{}a\allowbreak{}t\allowbreak{}e\allowbreak{}P\allowbreak{}r\allowbreak{}o\allowbreak{}g\allowbreak{}r\allowbreak{}e\allowbreak{}s\allowbreak{}s\allowbreak{}B\allowbreak{}a\allowbreak{}r\allowbreak{}(\allowbreak{})}
  \item
    \texttt{c\allowbreak{}r\allowbreak{}e\allowbreak{}a\allowbreak{}t\allowbreak{}e\allowbreak{}A\allowbreak{}n\allowbreak{}i\allowbreak{}m\allowbreak{}a\allowbreak{}t\allowbreak{}e\allowbreak{}d\allowbreak{}P\allowbreak{}r\allowbreak{}o\allowbreak{}g\allowbreak{}r\allowbreak{}e\allowbreak{}s\allowbreak{}s\allowbreak{}B\allowbreak{}a\allowbreak{}r\allowbreak{}(\allowbreak{})}
  \item
    \texttt{s\allowbreak{}p\allowbreak{}i\allowbreak{}n\allowbreak{}n\allowbreak{}e\allowbreak{}r\allowbreak{}F\allowbreak{}r\allowbreak{}a\allowbreak{}m\allowbreak{}e\allowbreak{}(\allowbreak{})}
  \item
    \texttt{c\allowbreak{}r\allowbreak{}e\allowbreak{}a\allowbreak{}t\allowbreak{}e\allowbreak{}S\allowbreak{}p\allowbreak{}i\allowbreak{}n\allowbreak{}n\allowbreak{}e\allowbreak{}r\allowbreak{}(\allowbreak{})}
  \end{itemize}
\item
  Variants:

  \begin{itemize}
  \tightlist
  \item
    determinate bar
  \item
    indeterminate spinner
  \item
    live controller helpers
  \item
    animated progress interpolation
  \end{itemize}
\item
  Use when:

  \begin{itemize}
  \tightlist
  \item
    the user needs honest feedback that work is ongoing
  \item
    determinate versus indeterminate progress is actually known
  \end{itemize}
\item
  Avoid when:

  \begin{itemize}
  \tightlist
  \item
    loading is so brief that the indicator would flicker more than it helps
  \item
    the state should be explained as a durable message or result instead of transient activity
  \end{itemize}
\item
  Content guidance:

  \begin{itemize}
  \tightlist
  \item
    use \texttt{p\allowbreak{}r\allowbreak{}o\allowbreak{}g\allowbreak{}r\allowbreak{}e\allowbreak{}s\allowbreak{}s\allowbreak{}B\allowbreak{}a\allowbreak{}r\allowbreak{}(\allowbreak{})} when completion can be estimated honestly
  \item
    use \texttt{s\allowbreak{}p\allowbreak{}i\allowbreak{}n\allowbreak{}n\allowbreak{}e\allowbreak{}r\allowbreak{}F\allowbreak{}r\allowbreak{}a\allowbreak{}m\allowbreak{}e\allowbreak{}(\allowbreak{})} / \texttt{c\allowbreak{}r\allowbreak{}e\allowbreak{}a\allowbreak{}t\allowbreak{}e\allowbreak{}S\allowbreak{}p\allowbreak{}i\allowbreak{}n\allowbreak{}n\allowbreak{}e\allowbreak{}r\allowbreak{}(\allowbreak{})} when activity is real but percent-complete is not
  \item
    switch from placeholder to progress indicator only when process feedback matters more than preserving layout shape
  \item
    pair progress indicators with clear labels so users know what is happening, not just that something is moving
  \end{itemize}
\item
  Ownership:

  \begin{itemize}
  \tightlist
  \item
    core
  \end{itemize}
\item
  Graceful lowering:

  \begin{itemize}
  \tightlist
  \item
    rich/static: preserve determinate or indeterminate activity cues honestly
  \item
    pipe: lower to explicit text progress or repeated status lines without fake animation semantics
  \item
    accessible: announce task state, progress, and completion explicitly in text
  \end{itemize}
\item
  Related families:

  \begin{itemize}
  \tightlist
  \item
    \texttt{s\allowbreak{}k\allowbreak{}e\allowbreak{}l\allowbreak{}e\allowbreak{}t\allowbreak{}o\allowbreak{}n\allowbreak{}(\allowbreak{})}
  \item
    \texttt{b\allowbreak{}a\allowbreak{}d\allowbreak{}g\allowbreak{}e\allowbreak{}(\allowbreak{})}
  \item
    notification system
  \end{itemize}
\item
  Carbon analogue:

  \begin{itemize}
  \tightlist
  \item
    progress bar / loading spinner
  \end{itemize}
\end{itemize}

\section{Structural families}\label{record-019-structural-families}

\subsection{Framed grouping}\label{record-019-framed-grouping}

\begin{itemize}
\tightlist
\item
  Family: \texttt{b\allowbreak{}o\allowbreak{}x\allowbreak{}(\allowbreak{})}
\item
  Variants:

  \begin{itemize}
  \tightlist
  \item
    titled via \texttt{h\allowbreak{}e\allowbreak{}a\allowbreak{}d\allowbreak{}e\allowbreak{}r\allowbreak{}B\allowbreak{}o\allowbreak{}x\allowbreak{}(\allowbreak{})}
  \item
    render-path companions via \texttt{b\allowbreak{}o\allowbreak{}x\allowbreak{}S\allowbreak{}u\allowbreak{}r\allowbreak{}f\allowbreak{}a\allowbreak{}c\allowbreak{}e\allowbreak{}(\allowbreak{})} and \texttt{h\allowbreak{}e\allowbreak{}a\allowbreak{}d\allowbreak{}e\allowbreak{}r\allowbreak{}B\allowbreak{}o\allowbreak{}x\allowbreak{}S\allowbreak{}u\allowbreak{}r\allowbreak{}f\allowbreak{}a\allowbreak{}c\allowbreak{}e\allowbreak{}(\allowbreak{})}
  \end{itemize}
\item
  Use when:

  \begin{itemize}
  \tightlist
  \item
    grouping, containment, or local region identity matter
  \item
    peer panels need to read as separate working areas
  \item
    a compact title plus detail helps orient a panel without consuming a full heading block
  \end{itemize}
\item
  Avoid when:

  \begin{itemize}
  \tightlist
  \item
    urgency or interruption is the primary job
  \item
    every subsection would get its own border purely for decoration
  \item
    a separator, heading, or whitespace would communicate structure more honestly
  \end{itemize}
\item
  Content guidance:

  \begin{itemize}
  \tightlist
  \item
    use \texttt{b\allowbreak{}o\allowbreak{}x\allowbreak{}(\allowbreak{})} for grouped content and \texttt{h\allowbreak{}e\allowbreak{}a\allowbreak{}d\allowbreak{}e\allowbreak{}r\allowbreak{}B\allowbreak{}o\allowbreak{}x\allowbreak{}(\allowbreak{})} when the region needs a clear title plus compact secondary detail
  \item
    titles should describe the region's job, not restate surrounding context
  \item
    \texttt{h\allowbreak{}e\allowbreak{}a\allowbreak{}d\allowbreak{}e\allowbreak{}r\allowbreak{}B\allowbreak{}o\allowbreak{}x\allowbreak{}(\allowbreak{})} detail text should carry terse metadata such as environment, path, version, or scope rather than a second sentence
  \item
    nested boxes should be rare and should explain real hierarchy or comparison, not just add more chrome
  \item
    keep dense prose out of narrow boxes unless wrapping materially helps comprehension
  \end{itemize}
\item
  Ownership:

  \begin{itemize}
  \tightlist
  \item
    core
  \end{itemize}
\item
  Graceful lowering:

  \begin{itemize}
  \tightlist
  \item
    rich/static: keep bordered or titled containment when space allows
  \item
    pipe: lower to plain grouped text with spacing and optional title
  \item
    accessible: preserve title and content order without decorative framing
  \end{itemize}
\item
  Related families:

  \begin{itemize}
  \tightlist
  \item
    \texttt{s\allowbreak{}e\allowbreak{}p\allowbreak{}a\allowbreak{}r\allowbreak{}a\allowbreak{}t\allowbreak{}o\allowbreak{}r\allowbreak{}(\allowbreak{})}
  \item
    \texttt{a\allowbreak{}l\allowbreak{}e\allowbreak{}r\allowbreak{}t\allowbreak{}(\allowbreak{})}
  \item
    \texttt{g\allowbreak{}r\allowbreak{}i\allowbreak{}d\allowbreak{}(\allowbreak{})}
  \item
    \texttt{s\allowbreak{}p\allowbreak{}l\allowbreak{}i\allowbreak{}t\allowbreak{}P\allowbreak{}a\allowbreak{}n\allowbreak{}e\allowbreak{}(\allowbreak{})}
  \end{itemize}
\item
  Carbon analogue:

  \begin{itemize}
  \tightlist
  \item
    tile / contained section
  \end{itemize}
\end{itemize}

\subsection{Explainability walkthroughs}\label{record-019-explainability-walkthroughs}

\begin{itemize}
\tightlist
\item
  Family: \texttt{e\allowbreak{}x\allowbreak{}p\allowbreak{}l\allowbreak{}a\allowbreak{}i\allowbreak{}n\allowbreak{}a\allowbreak{}b\allowbreak{}i\allowbreak{}l\allowbreak{}i\allowbreak{}t\allowbreak{}y\allowbreak{}(\allowbreak{})}
\item
  Variants:

  \begin{itemize}
  \tightlist
  \item
    visible provenance label
  \item
    rationale, evidence, next-action, and governance sections
  \item
    confidence and source metadata
  \end{itemize}
\item
  Use when:

  \begin{itemize}
  \tightlist
  \item
    AI-mediated or machine-assisted output needs explicit explanation instead of decorative summary
  \item
    the user needs one calm guided-flow surface that separates recommendation from supporting evidence
  \item
    the next action matters more than conversational prose
  \end{itemize}
\item
  Avoid when:

  \begin{itemize}
  \tightlist
  \item
    the content is just a generic status or note with no recommendation/evidence structure
  \item
    the app needs a non-AI guided flow or a full multi-step wizard instead of one explainability surface; prefer \texttt{g\allowbreak{}u\allowbreak{}i\allowbreak{}d\allowbreak{}e\allowbreak{}d\allowbreak{}F\allowbreak{}l\allowbreak{}o\allowbreak{}w\allowbreak{}(\allowbreak{})} for the calmer general case
  \item
    the app is hiding service latency or speculative output behind an authoritative-looking card
  \end{itemize}
\item
  Content guidance:

  \begin{itemize}
  \tightlist
  \item
    keep the title as the recommendation or explanation itself, not a vague marketing headline
  \item
    make \texttt{[\allowbreak{}A\allowbreak{}I\allowbreak{}]} or equivalent provenance visible
  \item
    separate rationale, evidence, and next action into distinct sections instead of mixing them together
  \item
    governance text should clarify review posture without becoming the dominant content
  \end{itemize}
\item
  Ownership:

  \begin{itemize}
  \tightlist
  \item
    core
  \end{itemize}
\item
  Graceful lowering:

  \begin{itemize}
  \tightlist
  \item
    rich/static: preserve section rhythm and explicit provenance inside one calm grouped surface
  \item
    pipe: lower to labeled plain-text sections with one obvious next action
  \item
    accessible: linearize the same fields explicitly with no dependence on borders, color, or layout
  \end{itemize}
\item
  Related families:

  \begin{itemize}
  \tightlist
  \item
    \texttt{g\allowbreak{}u\allowbreak{}i\allowbreak{}d\allowbreak{}e\allowbreak{}d\allowbreak{}F\allowbreak{}l\allowbreak{}o\allowbreak{}w\allowbreak{}(\allowbreak{})}
  \item
    \texttt{n\allowbreak{}o\allowbreak{}t\allowbreak{}e\allowbreak{}(\allowbreak{})}
  \item
    \texttt{a\allowbreak{}l\allowbreak{}e\allowbreak{}r\allowbreak{}t\allowbreak{}(\allowbreak{})}
  \item
    \texttt{s\allowbreak{}t\allowbreak{}e\allowbreak{}p\allowbreak{}p\allowbreak{}e\allowbreak{}r\allowbreak{}(\allowbreak{})}
  \item
    \texttt{t\allowbreak{}i\allowbreak{}m\allowbreak{}e\allowbreak{}l\allowbreak{}i\allowbreak{}n\allowbreak{}e\allowbreak{}(\allowbreak{})}
  \end{itemize}
\item
  Carbon analogue:

  \begin{itemize}
  \tightlist
  \item
    no exact analogue; closest to an explainability panel or guided recommendation card
  \end{itemize}
\end{itemize}

\subsection{Inspector panels}\label{record-019-inspector-panels}

\begin{itemize}
\tightlist
\item
  Family: \texttt{i\allowbreak{}n\allowbreak{}s\allowbreak{}p\allowbreak{}e\allowbreak{}c\allowbreak{}t\allowbreak{}o\allowbreak{}r\allowbreak{}(\allowbreak{})}
\item
  Variants:

  \begin{itemize}
  \tightlist
  \item
    current-selection emphasis
  \item
    compact titled sections
  \item
    muted supporting-copy sections
  \end{itemize}
\item
  Use when:

  \begin{itemize}
  \tightlist
  \item
    a side panel needs to summarize the currently selected thing without taking over the primary task
  \item
    supporting context should stay structured and calmer than the main content
  \item
    the panel needs one obvious current value plus compact supporting sections
  \end{itemize}
\item
  Avoid when:

  \begin{itemize}
  \tightlist
  \item
    the content is a full guided recommendation with evidence and next-action structure; prefer \texttt{e\allowbreak{}x\allowbreak{}p\allowbreak{}l\allowbreak{}a\allowbreak{}i\allowbreak{}n\allowbreak{}a\allowbreak{}b\allowbreak{}i\allowbreak{}l\allowbreak{}i\allowbreak{}t\allowbreak{}y\allowbreak{}(\allowbreak{})}
  \item
    the surface is really just an alert, note, or one-line status
  \item
    the panel needs its own complex navigation or multistep interaction model
  \end{itemize}
\item
  Content guidance:

  \begin{itemize}
  \tightlist
  \item
    make the current selection more obvious than supporting copy
  \item
    keep section titles short and concrete
  \item
    let supporting details look supporting; do not let descriptions overpower the active value
  \item
    stack explicit sections before inventing another decorative box inside the panel
  \end{itemize}
\item
  Ownership:

  \begin{itemize}
  \tightlist
  \item
    core
  \end{itemize}
\item
  Graceful lowering:

  \begin{itemize}
  \tightlist
  \item
    rich/static: preserve titled containment, current-selection emphasis, and compact section rhythm
  \item
    pipe: lower to explicit field labels with one obvious current selection
  \item
    accessible: linearize the same fields in plain language without borders or color
  \end{itemize}
\item
  Related families:

  \begin{itemize}
  \tightlist
  \item
    \texttt{b\allowbreak{}o\allowbreak{}x\allowbreak{}(\allowbreak{})}
  \item
    \texttt{e\allowbreak{}x\allowbreak{}p\allowbreak{}l\allowbreak{}a\allowbreak{}i\allowbreak{}n\allowbreak{}a\allowbreak{}b\allowbreak{}i\allowbreak{}l\allowbreak{}i\allowbreak{}t\allowbreak{}y\allowbreak{}(\allowbreak{})}
  \item
    \texttt{p\allowbreak{}r\allowbreak{}e\allowbreak{}f\allowbreak{}e\allowbreak{}r\allowbreak{}e\allowbreak{}n\allowbreak{}c\allowbreak{}e\allowbreak{}L\allowbreak{}i\allowbreak{}s\allowbreak{}t\allowbreak{}S\allowbreak{}u\allowbreak{}r\allowbreak{}f\allowbreak{}a\allowbreak{}c\allowbreak{}e\allowbreak{}(\allowbreak{})}
  \item
    notification system
  \item
    \texttt{i\allowbreak{}n\allowbreak{}s\allowbreak{}p\allowbreak{}e\allowbreak{}c\allowbreak{}t\allowbreak{}o\allowbreak{}r\allowbreak{}D\allowbreak{}r\allowbreak{}a\allowbreak{}w\allowbreak{}e\allowbreak{}r\allowbreak{}(\allowbreak{})}
  \end{itemize}
\item
  Carbon analogue:

  \begin{itemize}
  \tightlist
  \item
    side panel / contextual detail panel
  \end{itemize}
\end{itemize}

\subsection{Formatted documents and prose}\label{record-019-formatted-documents-and-prose}

\begin{itemize}
\tightlist
\item
  Family: \texttt{m\allowbreak{}a\allowbreak{}r\allowbreak{}k\allowbreak{}d\allowbreak{}o\allowbreak{}w\allowbreak{}n\allowbreak{}(\allowbreak{})}
\item
  Variants:

  \begin{itemize}
  \tightlist
  \item
    headings, emphasis, lists, blockquotes, links, GFM-style pipe tables, code fences, and other supported markdown structures with mode-aware lowering
  \end{itemize}
\item
  Use when:

  \begin{itemize}
  \tightlist
  \item
    help, reference, release notes, readme-like content, or bounded prose needs lightweight structure
  \item
    the same content should remain honest across rich, pipe, and accessible output modes
  \end{itemize}
\item
  Avoid when:

  \begin{itemize}
  \tightlist
  \item
    the content needs deep document navigation rather than one rendered block
  \item
    the app is really composing interface layout, forms, or command surfaces rather than prose
  \item
    browser-grade markdown fidelity or arbitrary user-authored rich documents are expected
  \end{itemize}
\item
  Content guidance:

  \begin{itemize}
  \tightlist
  \item
    headings should break the content into scannable chunks instead of reproducing the entire document hierarchy from another medium
  \item
    keep fenced code blocks and quotes short enough that the surrounding task still matters more than ornamental formatting
  \item
    links should remain explicit and self-describing, since terminals vary in clickable-link support
  \item
    use markdown for documents and reference prose, not as a substitute layout engine for application chrome
  \item
    when markdown appears in overlays or panes, keep the scope intentionally bounded and move to pager or navigation patterns if the content becomes a document-reading task
  \end{itemize}
\item
  Ownership:

  \begin{itemize}
  \tightlist
  \item
    core
  \end{itemize}
\item
  Graceful lowering:

  \begin{itemize}
  \tightlist
  \item
    rich/static: preserve supported markdown structure and emphasis honestly within terminal constraints
  \item
    pipe: lower to plain text with heading/list/code semantics still understandable without styling
  \item
    accessible: linearize headings, lists, links, and code cues explicitly in reading order
  \end{itemize}
\item
  Related families:

  \begin{itemize}
  \tightlist
  \item
    \texttt{n\allowbreak{}o\allowbreak{}t\allowbreak{}e\allowbreak{}(\allowbreak{})}
  \item
    \texttt{b\allowbreak{}o\allowbreak{}x\allowbreak{}(\allowbreak{})}
  \item
    \texttt{p\allowbreak{}a\allowbreak{}g\allowbreak{}e\allowbreak{}r\allowbreak{}(\allowbreak{})}
  \item
    \texttt{h\allowbreak{}e\allowbreak{}l\allowbreak{}p\allowbreak{}V\allowbreak{}i\allowbreak{}e\allowbreak{}w\allowbreak{}(\allowbreak{})}
  \end{itemize}
\item
  Carbon analogue:

  \begin{itemize}
  \tightlist
  \item
    structured text / markdown content block
  \end{itemize}
\end{itemize}

\subsection{Linked destinations}\label{record-019-linked-destinations}

\begin{itemize}
\tightlist
\item
  Family: \texttt{h\allowbreak{}y\allowbreak{}p\allowbreak{}e\allowbreak{}r\allowbreak{}l\allowbreak{}i\allowbreak{}n\allowbreak{}k\allowbreak{}(\allowbreak{})}
\item
  Variants:

  \begin{itemize}
  \tightlist
  \item
    OSC 8 clickable link
  \item
    explicit fallback text or URL modes
  \end{itemize}
\item
  Use when:

  \begin{itemize}
  \tightlist
  \item
    the destination matters and should remain part of the rendered output
  \item
    terminals with hyperlink support should get a direct clickable affordance without hiding the actual destination semantics
  \end{itemize}
\item
  Avoid when:

  \begin{itemize}
  \tightlist
  \item
    the destination is untrusted, ambiguous, or should not be activated casually
  \item
    link text is the only place where critical meaning or safety context appears
  \item
    the user needs an app-owned action rather than an external destination
  \end{itemize}
\item
  Content guidance:

  \begin{itemize}
  \tightlist
  \item
    link labels should describe the destination or resource, not generic phrases like ``click here''
  \item
    when trust matters, keep the host or destination obvious in the surrounding text or fallback behavior
  \item
    use fallback modes intentionally so non-supporting terminals still preserve the destination honestly
  \end{itemize}
\item
  Ownership:

  \begin{itemize}
  \tightlist
  \item
    core
  \end{itemize}
\item
  Graceful lowering:

  \begin{itemize}
  \tightlist
  \item
    rich/static: use OSC 8 when supported while keeping meaningful visible link text
  \item
    pipe: lower to explicit text plus destination URL when needed
  \item
    accessible: preserve label and destination clearly in reading order without assuming clickability
  \end{itemize}
\item
  Related families:

  \begin{itemize}
  \tightlist
  \item
    \texttt{m\allowbreak{}a\allowbreak{}r\allowbreak{}k\allowbreak{}d\allowbreak{}o\allowbreak{}w\allowbreak{}n\allowbreak{}(\allowbreak{})}
  \item
    \texttt{n\allowbreak{}o\allowbreak{}t\allowbreak{}e\allowbreak{}(\allowbreak{})}
  \item
    \texttt{h\allowbreak{}e\allowbreak{}l\allowbreak{}p\allowbreak{}V\allowbreak{}i\allowbreak{}e\allowbreak{}w\allowbreak{}(\allowbreak{})}
  \end{itemize}
\item
  Carbon analogue:

  \begin{itemize}
  \tightlist
  \item
    link
  \end{itemize}
\end{itemize}

\subsection{Dividers}\label{record-019-dividers}

\begin{itemize}
\tightlist
\item
  Family: \texttt{s\allowbreak{}e\allowbreak{}p\allowbreak{}a\allowbreak{}r\allowbreak{}a\allowbreak{}t\allowbreak{}o\allowbreak{}r\allowbreak{}(\allowbreak{})}
\item
  Variants:

  \begin{itemize}
  \tightlist
  \item
    labeled, unlabeled
  \item
    render-path companion: \texttt{s\allowbreak{}e\allowbreak{}p\allowbreak{}a\allowbreak{}r\allowbreak{}a\allowbreak{}t\allowbreak{}o\allowbreak{}r\allowbreak{}S\allowbreak{}u\allowbreak{}r\allowbreak{}f\allowbreak{}a\allowbreak{}c\allowbreak{}e\allowbreak{}(\allowbreak{})}
  \end{itemize}
\item
  Use when:

  \begin{itemize}
  \tightlist
  \item
    a section boundary is needed without full containment
  \end{itemize}
\item
  Avoid when:

  \begin{itemize}
  \tightlist
  \item
    the content needs its own grouped region
  \end{itemize}
\item
  Content guidance:

  \begin{itemize}
  \tightlist
  \item
    labels should name the next section or state, not repeat a page title
  \item
    use dividers sparingly so they mark real boundaries rather than visual noise
  \end{itemize}
\item
  Ownership:

  \begin{itemize}
  \tightlist
  \item
    core
  \end{itemize}
\item
  Graceful lowering:

  \begin{itemize}
  \tightlist
  \item
    rich/static: render as visual divider or labeled rule
  \item
    pipe: lower to simple text divider or heading-like label
  \item
    accessible: preserve section naming without decorative characters being the only cue
  \end{itemize}
\item
  Related families:

  \begin{itemize}
  \tightlist
  \item
    \texttt{b\allowbreak{}o\allowbreak{}x\allowbreak{}(\allowbreak{})}
  \item
    \texttt{t\allowbreak{}a\allowbreak{}b\allowbreak{}s\allowbreak{}(\allowbreak{})}
  \item
    \texttt{b\allowbreak{}r\allowbreak{}e\allowbreak{}a\allowbreak{}d\allowbreak{}c\allowbreak{}r\allowbreak{}u\allowbreak{}m\allowbreak{}b\allowbreak{}(\allowbreak{})}
  \end{itemize}
\item
  Carbon analogue:

  \begin{itemize}
  \tightlist
  \item
    divider
  \end{itemize}
\end{itemize}

\section{Choice and form families}\label{record-019-choice-and-form-families}

\subsection{Text entry}\label{record-019-text-entry}

\begin{itemize}
\tightlist
\item
  Family:

  \begin{itemize}
  \tightlist
  \item
    \texttt{i\allowbreak{}n\allowbreak{}p\allowbreak{}u\allowbreak{}t\allowbreak{}(\allowbreak{})}
  \item
    \texttt{t\allowbreak{}e\allowbreak{}x\allowbreak{}t\allowbreak{}a\allowbreak{}r\allowbreak{}e\allowbreak{}a\allowbreak{}(\allowbreak{})}
  \end{itemize}
\item
  Variants:

  \begin{itemize}
  \tightlist
  \item
    short-form versus long-form entry
  \end{itemize}
\item
  Use when:

  \begin{itemize}
  \tightlist
  \item
    the user is providing text rather than choosing from a fixed set
  \end{itemize}
\item
  Avoid when:

  \begin{itemize}
  \tightlist
  \item
    choices are already known and enumerable
  \end{itemize}
\item
  Content guidance:

  \begin{itemize}
  \tightlist
  \item
    prompts should describe the expected input shape and any critical constraints
  \item
    placeholders and defaults should clarify intent, not substitute for the label
  \end{itemize}
\item
  Ownership:

  \begin{itemize}
  \tightlist
  \item
    core
  \end{itemize}
\item
  Graceful lowering:

  \begin{itemize}
  \tightlist
  \item
    rich: interactive field/editor treatment with validation feedback
  \item
    static: render current value or prompt snapshot honestly
  \item
    pipe: fall back to line-buffered prompt/input flow
  \item
    accessible: keep prompt, validation, and submitted value explicit in text
  \end{itemize}
\item
  Related families:

  \begin{itemize}
  \tightlist
  \item
    \texttt{s\allowbreak{}e\allowbreak{}l\allowbreak{}e\allowbreak{}c\allowbreak{}t\allowbreak{}(\allowbreak{})}
  \item
    \texttt{f\allowbreak{}i\allowbreak{}l\allowbreak{}t\allowbreak{}e\allowbreak{}r\allowbreak{}(\allowbreak{})}
  \item
    \texttt{g\allowbreak{}r\allowbreak{}o\allowbreak{}u\allowbreak{}p\allowbreak{}(\allowbreak{})}
  \item
    \texttt{w\allowbreak{}i\allowbreak{}z\allowbreak{}a\allowbreak{}r\allowbreak{}d\allowbreak{}(\allowbreak{})}
  \end{itemize}
\item
  Carbon analogue:

  \begin{itemize}
  \tightlist
  \item
    text input / text area
  \end{itemize}
\end{itemize}

\subsection{Single choice}\label{record-019-single-choice}

\begin{itemize}
\tightlist
\item
  Family:

  \begin{itemize}
  \tightlist
  \item
    \texttt{s\allowbreak{}e\allowbreak{}l\allowbreak{}e\allowbreak{}c\allowbreak{}t\allowbreak{}(\allowbreak{})}
  \item
    \texttt{f\allowbreak{}i\allowbreak{}l\allowbreak{}t\allowbreak{}e\allowbreak{}r\allowbreak{}(\allowbreak{})}
  \end{itemize}
\item
  Variants:

  \begin{itemize}
  \tightlist
  \item
    regular choice
  \item
    searchable choice
  \end{itemize}
\item
  Use when:

  \begin{itemize}
  \tightlist
  \item
    the result is stored value state
  \end{itemize}
\item
  Avoid when:

  \begin{itemize}
  \tightlist
  \item
    selecting should trigger an app command rather than persist a value
  \item
    the option set is large enough that search, ranking, or narrowing is the real task
  \end{itemize}
\item
  Content guidance:

  \begin{itemize}
  \tightlist
  \item
    option labels should be distinct, parallel, and easy to scan
  \item
    searchable choice should use concise matching text, not long descriptive paragraphs
  \item
    if search is the real job, prefer \texttt{f\allowbreak{}i\allowbreak{}l\allowbreak{}t\allowbreak{}e\allowbreak{}r\allowbreak{}(\allowbreak{})} semantics over stuffing too many options into plain \texttt{s\allowbreak{}e\allowbreak{}l\allowbreak{}e\allowbreak{}c\allowbreak{}t\allowbreak{}(\allowbreak{})}
  \end{itemize}
\item
  Ownership:

  \begin{itemize}
  \tightlist
  \item
    core
  \end{itemize}
\item
  Graceful lowering:

  \begin{itemize}
  \tightlist
  \item
    rich: keyboard choice list with focus and optional search narrowing
  \item
    static: show current selection or available options snapshot
  \item
    pipe: lower to numbered or searchable textual selection flow
  \item
    accessible: keep option labels, current selection, and prompt semantics explicit
  \end{itemize}
\item
  Related families:

  \begin{itemize}
  \tightlist
  \item
    \texttt{m\allowbreak{}u\allowbreak{}l\allowbreak{}t\allowbreak{}i\allowbreak{}s\allowbreak{}e\allowbreak{}l\allowbreak{}e\allowbreak{}c\allowbreak{}t\allowbreak{}(\allowbreak{})}
  \item
    \texttt{c\allowbreak{}o\allowbreak{}m\allowbreak{}m\allowbreak{}a\allowbreak{}n\allowbreak{}d\allowbreak{}P\allowbreak{}a\allowbreak{}l\allowbreak{}e\allowbreak{}t\allowbreak{}t\allowbreak{}e\allowbreak{}(\allowbreak{})}
  \item
    \texttt{i\allowbreak{}n\allowbreak{}p\allowbreak{}u\allowbreak{}t\allowbreak{}(\allowbreak{})}
  \end{itemize}
\item
  Carbon analogue:

  \begin{itemize}
  \tightlist
  \item
    dropdown / combo box
  \end{itemize}
\end{itemize}

\subsection{Multiple choice}\label{record-019-multiple-choice}

\begin{itemize}
\tightlist
\item
  Family: \texttt{m\allowbreak{}u\allowbreak{}l\allowbreak{}t\allowbreak{}i\allowbreak{}s\allowbreak{}e\allowbreak{}l\allowbreak{}e\allowbreak{}c\allowbreak{}t\allowbreak{}(\allowbreak{})}
\item
  Variants:

  \begin{itemize}
  \tightlist
  \item
    default selections, checkbox set behavior
  \end{itemize}
\item
  Use when:

  \begin{itemize}
  \tightlist
  \item
    the user is building a set
  \end{itemize}
\item
  Avoid when:

  \begin{itemize}
  \tightlist
  \item
    the choice is singular or immediate action-oriented
  \item
    the user is really firing commands rather than collecting lasting state
  \end{itemize}
\item
  Content guidance:

  \begin{itemize}
  \tightlist
  \item
    options should read like members of one set, not unrelated commands
  \item
    summary text should make selection state obvious without requiring visual checkboxes
  \end{itemize}
\item
  Ownership:

  \begin{itemize}
  \tightlist
  \item
    core
  \end{itemize}
\item
  Graceful lowering:

  \begin{itemize}
  \tightlist
  \item
    rich: checkbox-style set selection with keyboard toggling
  \item
    static: show selected values snapshot honestly
  \item
    pipe: lower to comma-separated or numbered textual selection flow
  \item
    accessible: describe the current set in plain text rather than relying on checkmarks alone
  \end{itemize}
\item
  Related families:

  \begin{itemize}
  \tightlist
  \item
    \texttt{s\allowbreak{}e\allowbreak{}l\allowbreak{}e\allowbreak{}c\allowbreak{}t\allowbreak{}(\allowbreak{})}
  \item
    \texttt{f\allowbreak{}i\allowbreak{}l\allowbreak{}t\allowbreak{}e\allowbreak{}r\allowbreak{}(\allowbreak{})}
  \item
    \texttt{g\allowbreak{}r\allowbreak{}o\allowbreak{}u\allowbreak{}p\allowbreak{}(\allowbreak{})}
  \end{itemize}
\item
  Carbon analogue:

  \begin{itemize}
  \tightlist
  \item
    multi-select / checkbox group
  \end{itemize}
\end{itemize}

\subsection{Binary decision}\label{record-019-binary-decision}

\begin{itemize}
\tightlist
\item
  Family: \texttt{c\allowbreak{}o\allowbreak{}n\allowbreak{}f\allowbreak{}i\allowbreak{}r\allowbreak{}m\allowbreak{}(\allowbreak{})}
\item
  Variants:

  \begin{itemize}
  \tightlist
  \item
    default yes/no orientation
  \end{itemize}
\item
  Use when:

  \begin{itemize}
  \tightlist
  \item
    the decision is genuinely binary
  \end{itemize}
\item
  Avoid when:

  \begin{itemize}
  \tightlist
  \item
    nuance, comparison, or explanation is needed
  \item
    the prompt is really a multi-option choice disguised as yes/no
  \end{itemize}
\item
  Content guidance:

  \begin{itemize}
  \tightlist
  \item
    questions should be explicit about the consequence of yes versus no
  \item
    destructive confirmations should name the thing being changed or removed
  \end{itemize}
\item
  Ownership:

  \begin{itemize}
  \tightlist
  \item
    core
  \end{itemize}
\item
  Graceful lowering:

  \begin{itemize}
  \tightlist
  \item
    rich: yes/no confirmation prompt or focused confirm surface
  \item
    static: show the confirmation question and current default honestly
  \item
    pipe: lower naturally to textual yes/no input
  \item
    accessible: preserve the binary choice and default state explicitly
  \end{itemize}
\item
  Related families:

  \begin{itemize}
  \tightlist
  \item
    \texttt{m\allowbreak{}o\allowbreak{}d\allowbreak{}a\allowbreak{}l\allowbreak{}(\allowbreak{})}
  \item
    \texttt{a\allowbreak{}l\allowbreak{}e\allowbreak{}r\allowbreak{}t\allowbreak{}(\allowbreak{})}
  \item
    \texttt{w\allowbreak{}i\allowbreak{}z\allowbreak{}a\allowbreak{}r\allowbreak{}d\allowbreak{}(\allowbreak{})}
  \end{itemize}
\item
  Carbon analogue:

  \begin{itemize}
  \tightlist
  \item
    confirmation dialog pattern
  \end{itemize}
\end{itemize}

\subsection{Multi-field and staged forms}\label{record-019-multi-field-and-staged-forms}

\begin{itemize}
\tightlist
\item
  Family:

  \begin{itemize}
  \tightlist
  \item
    \texttt{g\allowbreak{}r\allowbreak{}o\allowbreak{}u\allowbreak{}p\allowbreak{}(\allowbreak{})}
  \item
    \texttt{w\allowbreak{}i\allowbreak{}z\allowbreak{}a\allowbreak{}r\allowbreak{}d\allowbreak{}(\allowbreak{})}
  \end{itemize}
\item
  Variants:

  \begin{itemize}
  \tightlist
  \item
    single-step grouped form
  \item
    staged or branching flow
  \end{itemize}
\item
  Use when:

  \begin{itemize}
  \tightlist
  \item
    the user is progressing through related inputs
  \end{itemize}
\item
  Avoid when:

  \begin{itemize}
  \tightlist
  \item
    the fields are unrelated or the result is a command
  \end{itemize}
\item
  Content guidance:

  \begin{itemize}
  \tightlist
  \item
    each step or group should have a clear goal and avoid mixing unrelated concepts
  \item
    summaries and progress text should orient the user without duplicating every field label
  \end{itemize}
\item
  Ownership:

  \begin{itemize}
  \tightlist
  \item
    core
  \end{itemize}
\item
  Graceful lowering:

  \begin{itemize}
  \tightlist
  \item
    rich: grouped or staged prompts with progress and validation
  \item
    static: show the current step/group snapshot honestly
  \item
    pipe: lower to sequential prompts while preserving branching meaning
  \item
    accessible: keep step names, requirements, and progress explicit in text
  \end{itemize}
\item
  Related families:

  \begin{itemize}
  \tightlist
  \item
    \texttt{i\allowbreak{}n\allowbreak{}p\allowbreak{}u\allowbreak{}t\allowbreak{}(\allowbreak{})}
  \item
    \texttt{s\allowbreak{}e\allowbreak{}l\allowbreak{}e\allowbreak{}c\allowbreak{}t\allowbreak{}(\allowbreak{})}
  \item
    \texttt{c\allowbreak{}o\allowbreak{}n\allowbreak{}f\allowbreak{}i\allowbreak{}r\allowbreak{}m\allowbreak{}(\allowbreak{})}
  \item
    \texttt{s\allowbreak{}t\allowbreak{}e\allowbreak{}p\allowbreak{}p\allowbreak{}e\allowbreak{}r\allowbreak{}(\allowbreak{})}
  \end{itemize}
\item
  Carbon analogue:

  \begin{itemize}
  \tightlist
  \item
    form group / multi-step flow
  \end{itemize}
\end{itemize}

\section{Content organization families}\label{record-019-content-organization-families}

\subsection{Peer navigation}\label{record-019-peer-navigation}

\begin{itemize}
\tightlist
\item
  Family: \texttt{t\allowbreak{}a\allowbreak{}b\allowbreak{}s\allowbreak{}(\allowbreak{})}
\item
  Variants:

  \begin{itemize}
  \tightlist
  \item
    tab labels, badges, active state
  \end{itemize}
\item
  Use when:

  \begin{itemize}
  \tightlist
  \item
    views are peers and only one is active
  \end{itemize}
\item
  Avoid when:

  \begin{itemize}
  \tightlist
  \item
    order or hierarchy is the main idea
  \end{itemize}
\item
  Content guidance:

  \begin{itemize}
  \tightlist
  \item
    tab labels should be short destination names, not action phrases
  \item
    unread counts or compact status belong in secondary adornments, not the core label
  \end{itemize}
\item
  Ownership:

  \begin{itemize}
  \tightlist
  \item
    core
  \end{itemize}
\item
  Graceful lowering:

  \begin{itemize}
  \tightlist
  \item
    rich/static: keep one active peer view visibly selected
  \item
    pipe: lower to labeled current section plus sibling list if needed
  \item
    accessible: preserve peer destinations and active state explicitly in text
  \end{itemize}
\item
  Related families:

  \begin{itemize}
  \tightlist
  \item
    \texttt{s\allowbreak{}t\allowbreak{}e\allowbreak{}p\allowbreak{}p\allowbreak{}e\allowbreak{}r\allowbreak{}(\allowbreak{})}
  \item
    \texttt{b\allowbreak{}r\allowbreak{}e\allowbreak{}a\allowbreak{}d\allowbreak{}c\allowbreak{}r\allowbreak{}u\allowbreak{}m\allowbreak{}b\allowbreak{}(\allowbreak{})}
  \item
    \texttt{c\allowbreak{}r\allowbreak{}e\allowbreak{}a\allowbreak{}t\allowbreak{}e\allowbreak{}F\allowbreak{}r\allowbreak{}a\allowbreak{}m\allowbreak{}e\allowbreak{}d\allowbreak{}A\allowbreak{}p\allowbreak{}p\allowbreak{}(\allowbreak{})}
  \end{itemize}
\item
  Carbon analogue:

  \begin{itemize}
  \tightlist
  \item
    tabs
  \end{itemize}
\end{itemize}

\subsection{Progressive disclosure}\label{record-019-progressive-disclosure}

\begin{itemize}
\tightlist
\item
  Family:

  \begin{itemize}
  \tightlist
  \item
    \texttt{a\allowbreak{}c\allowbreak{}c\allowbreak{}o\allowbreak{}r\allowbreak{}d\allowbreak{}i\allowbreak{}o\allowbreak{}n\allowbreak{}(\allowbreak{})}
  \item
    \texttt{i\allowbreak{}n\allowbreak{}t\allowbreak{}e\allowbreak{}r\allowbreak{}a\allowbreak{}c\allowbreak{}t\allowbreak{}i\allowbreak{}v\allowbreak{}e\allowbreak{}A\allowbreak{}c\allowbreak{}c\allowbreak{}o\allowbreak{}r\allowbreak{}d\allowbreak{}i\allowbreak{}o\allowbreak{}n\allowbreak{}(\allowbreak{})}
  \end{itemize}
\item
  Variants:

  \begin{itemize}
  \tightlist
  \item
    static disclosure
  \item
    keyboard-owned TUI inspection
  \end{itemize}
\item
  Use when:

  \begin{itemize}
  \tightlist
  \item
    detail is secondary to summary
  \end{itemize}
\item
  Avoid when:

  \begin{itemize}
  \tightlist
  \item
    sections are peers that deserve first-class navigation
  \end{itemize}
\item
  Content guidance:

  \begin{itemize}
  \tightlist
  \item
    section headers should summarize the concealed content well enough to support scanning
  \item
    disclosed content should stay tightly related to its summary row
  \end{itemize}
\item
  Ownership:

  \begin{itemize}
  \tightlist
  \item
    core plus TUI interaction layer
  \end{itemize}
\item
  Graceful lowering:

  \begin{itemize}
  \tightlist
  \item
    rich: collapsible sections with optional keyboard-owned inspection
  \item
    static: show expanded/collapsed state honestly without requiring motion
  \item
    pipe: lower to section headings with visible disclosed content
  \item
    accessible: keep section labels and disclosure state explicit in text
  \end{itemize}
\item
  Related families:

  \begin{itemize}
  \tightlist
  \item
    \texttt{t\allowbreak{}a\allowbreak{}b\allowbreak{}s\allowbreak{}(\allowbreak{})}
  \item
    \texttt{b\allowbreak{}o\allowbreak{}x\allowbreak{}(\allowbreak{})}
  \item
    \texttt{i\allowbreak{}n\allowbreak{}t\allowbreak{}e\allowbreak{}r\allowbreak{}a\allowbreak{}c\allowbreak{}t\allowbreak{}i\allowbreak{}v\allowbreak{}e\allowbreak{}A\allowbreak{}c\allowbreak{}c\allowbreak{}o\allowbreak{}r\allowbreak{}d\allowbreak{}i\allowbreak{}o\allowbreak{}n\allowbreak{}(\allowbreak{})}
  \end{itemize}
\item
  Carbon analogue:

  \begin{itemize}
  \tightlist
  \item
    accordion
  \end{itemize}
\end{itemize}

\subsection{Path and progress}\label{record-019-path-and-progress}

\begin{itemize}
\tightlist
\item
  Family:

  \begin{itemize}
  \tightlist
  \item
    \texttt{b\allowbreak{}r\allowbreak{}e\allowbreak{}a\allowbreak{}d\allowbreak{}c\allowbreak{}r\allowbreak{}u\allowbreak{}m\allowbreak{}b\allowbreak{}(\allowbreak{})}
  \item
    \texttt{s\allowbreak{}t\allowbreak{}e\allowbreak{}p\allowbreak{}p\allowbreak{}e\allowbreak{}r\allowbreak{}(\allowbreak{})}
  \item
    \texttt{p\allowbreak{}a\allowbreak{}g\allowbreak{}i\allowbreak{}n\allowbreak{}a\allowbreak{}t\allowbreak{}o\allowbreak{}r\allowbreak{}(\allowbreak{})}
  \end{itemize}
\item
  Variants:

  \begin{itemize}
  \tightlist
  \item
    location, process progress, page count
  \end{itemize}
\item
  Use when:

  \begin{itemize}
  \tightlist
  \item
    the UI needs wayfinding or progress communication
  \end{itemize}
\item
  Avoid when:

  \begin{itemize}
  \tightlist
  \item
    tabs or lists express the job more clearly
  \end{itemize}
\item
  Content guidance:

  \begin{itemize}
  \tightlist
  \item
    path and progress labels should emphasize where the user is, not every possible destination
  \item
    step titles should be stable and action-oriented enough to support review and recovery
  \item
    use \texttt{b\allowbreak{}r\allowbreak{}e\allowbreak{}a\allowbreak{}d\allowbreak{}c\allowbreak{}r\allowbreak{}u\allowbreak{}m\allowbreak{}b\allowbreak{}(\allowbreak{})} when path context helps orient the current location, not as a decorative header flourish
  \item
    use \texttt{p\allowbreak{}a\allowbreak{}g\allowbreak{}i\allowbreak{}n\allowbreak{}a\allowbreak{}t\allowbreak{}o\allowbreak{}r\allowbreak{}(\allowbreak{})} when compact position-in-sequence feedback is enough and the user does not need a richer navigation control
  \item
    use \texttt{s\allowbreak{}t\allowbreak{}e\allowbreak{}p\allowbreak{}p\allowbreak{}e\allowbreak{}r\allowbreak{}(\allowbreak{})} when the current stage and remaining path matter more than peer switching
  \end{itemize}
\item
  Ownership:

  \begin{itemize}
  \tightlist
  \item
    core
  \end{itemize}
\item
  Graceful lowering:

  \begin{itemize}
  \tightlist
  \item
    rich/static: retain location, step, or page state with simple styling
  \item
    pipe: lower to plain path/progress/page summaries
  \item
    accessible: preserve order and active/current state explicitly
  \end{itemize}
\item
  Related families:

  \begin{itemize}
  \tightlist
  \item
    \texttt{t\allowbreak{}a\allowbreak{}b\allowbreak{}s\allowbreak{}(\allowbreak{})}
  \item
    \texttt{w\allowbreak{}i\allowbreak{}z\allowbreak{}a\allowbreak{}r\allowbreak{}d\allowbreak{}(\allowbreak{})}
  \item
    \texttt{s\allowbreak{}t\allowbreak{}a\allowbreak{}t\allowbreak{}u\allowbreak{}s\allowbreak{}B\allowbreak{}a\allowbreak{}r\allowbreak{}(\allowbreak{})}
  \end{itemize}
\item
  Carbon analogue:

  \begin{itemize}
  \tightlist
  \item
    breadcrumb / progress indicator / pagination
  \end{itemize}
\end{itemize}

\subsection{Motion and shader effects}\label{record-019-motion-and-shader-effects}

\begin{itemize}
\tightlist
\item
  Family:

  \begin{itemize}
  \tightlist
  \item
    \texttt{c\allowbreak{}a\allowbreak{}n\allowbreak{}v\allowbreak{}a\allowbreak{}s\allowbreak{}(\allowbreak{})}
  \item
    transition shaders
  \item
    \texttt{a\allowbreak{}n\allowbreak{}i\allowbreak{}m\allowbreak{}a\allowbreak{}t\allowbreak{}e\allowbreak{}(\allowbreak{})}
  \item
    \texttt{t\allowbreak{}i\allowbreak{}m\allowbreak{}e\allowbreak{}l\allowbreak{}i\allowbreak{}n\allowbreak{}e\allowbreak{}(\allowbreak{})}
  \end{itemize}
\item
  Variants:

  \begin{itemize}
  \tightlist
  \item
    decorative visual field
  \item
    transition reinforcement
  \item
    motion choreography
  \item
    custom shader override
  \end{itemize}
\item
  Use when:

  \begin{itemize}
  \tightlist
  \item
    motion or shader output clarifies state change, transition, or atmosphere
  \end{itemize}
\item
  Avoid when:

  \begin{itemize}
  \tightlist
  \item
    the effect is only decorative noise
  \item
    readability, scanning, or task focus would be harmed
  \item
    the same meaning would be clearer as stable content or a simpler status cue
  \end{itemize}
\item
  Content guidance:

  \begin{itemize}
  \tightlist
  \item
    \texttt{c\allowbreak{}a\allowbreak{}n\allowbreak{}v\allowbreak{}a\allowbreak{}s\allowbreak{}(\allowbreak{})} should be reserved for deliberate visual moments such as splash surfaces, atmospheric backgrounds, or specialized visual regions, not routine productivity chrome
  \item
    transition shaders should reinforce page or workspace change, not compete with the content being changed
  \item
    animation timing should support comprehension and must have an honest reduced-motion or non-interactive fallback
  \item
    shader overrides should preserve semantic content where possible and avoid turning ordinary task flows into spectacle
  \end{itemize}
\item
  Ownership:

  \begin{itemize}
  \tightlist
  \item
    TUI
  \end{itemize}
\item
  Graceful lowering:

  \begin{itemize}
  \tightlist
  \item
    rich: preserve the motion or visual effect when it materially helps
  \item
    static: lower to truthful final-state snapshots without pretending the transition still exists
  \item
    pipe: drop decorative effects and keep only the meaning-bearing content
  \item
    accessible: preserve state change meaning explicitly without requiring visual motion
  \end{itemize}
\item
  Related families:

  \begin{itemize}
  \tightlist
  \item
    \texttt{c\allowbreak{}r\allowbreak{}e\allowbreak{}a\allowbreak{}t\allowbreak{}e\allowbreak{}F\allowbreak{}r\allowbreak{}a\allowbreak{}m\allowbreak{}e\allowbreak{}d\allowbreak{}A\allowbreak{}p\allowbreak{}p\allowbreak{}(\allowbreak{})}
  \item
    \texttt{s\allowbreak{}t\allowbreak{}a\allowbreak{}t\allowbreak{}u\allowbreak{}s\allowbreak{}B\allowbreak{}a\allowbreak{}r\allowbreak{}(\allowbreak{})}
  \item
    notification system
  \item
    \texttt{t\allowbreak{}i\allowbreak{}m\allowbreak{}e\allowbreak{}l\allowbreak{}i\allowbreak{}n\allowbreak{}e\allowbreak{}(\allowbreak{})}
  \end{itemize}
\item
  Carbon analogue:

  \begin{itemize}
  \tightlist
  \item
    no exact analogue; closest to expressive motion and illustration guidance rather than one component family
  \end{itemize}
\end{itemize}

\subsection{Expressive branding and decorative emphasis}\label{record-019-expressive-branding-and-decorative-emphasis}

\begin{itemize}
\tightlist
\item
  Family:

  \begin{itemize}
  \tightlist
  \item
    \texttt{l\allowbreak{}o\allowbreak{}a\allowbreak{}d\allowbreak{}R\allowbreak{}a\allowbreak{}n\allowbreak{}d\allowbreak{}o\allowbreak{}m\allowbreak{}L\allowbreak{}o\allowbreak{}g\allowbreak{}o\allowbreak{}(\allowbreak{})}
  \item
    \texttt{g\allowbreak{}r\allowbreak{}a\allowbreak{}d\allowbreak{}i\allowbreak{}e\allowbreak{}n\allowbreak{}t\allowbreak{}T\allowbreak{}e\allowbreak{}x\allowbreak{}t\allowbreak{}(\allowbreak{})}
  \end{itemize}
\item
  Variants:

  \begin{itemize}
  \tightlist
  \item
    ASCII logo sizes
  \item
    single-line or multiline gradient emphasis
  \end{itemize}
\item
  Use when:

  \begin{itemize}
  \tightlist
  \item
    the interface needs a deliberate branded or celebratory moment
  \item
    splash surfaces, docs, demos, and launch moments benefit from expressive treatment
  \end{itemize}
\item
  Avoid when:

  \begin{itemize}
  \tightlist
  \item
    routine application chrome or task-critical labels need to stay plain and scannable
  \item
    color, decoration, or branding would compete with the actual work
  \item
    the same text must remain equally legible in constrained or no-color environments
  \end{itemize}
\item
  Content guidance:

  \begin{itemize}
  \tightlist
  \item
    treat logos and gradients as rare emphasis, not the default voice of the interface
  \item
    never make critical state, instructions, or navigation depend on gradient color alone
  \item
    use expressive branding to open, celebrate, or orient, then get out of the user's way
  \end{itemize}
\item
  Ownership:

  \begin{itemize}
  \tightlist
  \item
    core
  \end{itemize}
\item
  Graceful lowering:

  \begin{itemize}
  \tightlist
  \item
    rich/static: preserve expressive branding where it adds real atmosphere without obscuring content
  \item
    pipe: lower to plain ASCII/text without pretending the gradient or brand treatment still carries meaning
  \item
    accessible: keep the textual content explicit and avoid decorative output that adds noise to the reading order
  \end{itemize}
\item
  Related families:

  \begin{itemize}
  \tightlist
  \item
    \texttt{c\allowbreak{}a\allowbreak{}n\allowbreak{}v\allowbreak{}a\allowbreak{}s\allowbreak{}(\allowbreak{})}
  \item
    \texttt{b\allowbreak{}o\allowbreak{}x\allowbreak{}(\allowbreak{})}
  \item
    \texttt{m\allowbreak{}a\allowbreak{}r\allowbreak{}k\allowbreak{}d\allowbreak{}o\allowbreak{}w\allowbreak{}n\allowbreak{}(\allowbreak{})}
  \end{itemize}
\item
  Carbon analogue:

  \begin{itemize}
  \tightlist
  \item
    closest to illustration, expressive color, and brand guidance rather than one component family
  \end{itemize}
\end{itemize}

\section{Authoring families}\label{record-019-authoring-families}

\subsection{Mode-aware custom primitives}\label{record-019-mode-aware-custom-primitives}

\begin{itemize}
\tightlist
\item
  Family:

  \begin{itemize}
  \tightlist
  \item
    \texttt{r\allowbreak{}e\allowbreak{}n\allowbreak{}d\allowbreak{}e\allowbreak{}r\allowbreak{}B\allowbreak{}y\allowbreak{}M\allowbreak{}o\allowbreak{}d\allowbreak{}e\allowbreak{}(\allowbreak{})}
  \end{itemize}
\item
  Variants:

  \begin{itemize}
  \tightlist
  \item
    interactive/rich branch
  \item
    pipe branch
  \item
    accessible branch
  \item
    app-specific authored primitive
  \end{itemize}
\item
  Use when:

  \begin{itemize}
  \tightlist
  \item
    the app needs a domain-specific primitive that does not belong in the shared component library
  \item
    the same semantic thing must render honestly across multiple output modes
  \end{itemize}
\item
  Avoid when:

  \begin{itemize}
  \tightlist
  \item
    an existing Bijou family already matches the job
  \item
    the branching only exists to chase cosmetics rather than preserve semantic truth
  \end{itemize}
\item
  Content guidance:

  \begin{itemize}
  \tightlist
  \item
    author one semantic primitive first, then define how it lowers by mode
  \item
    keep all branches truthful to the same underlying meaning instead of inventing different features per mode
  \item
    prefer app-owned names and domain language over generic helper wrappers
  \end{itemize}
\item
  Ownership:

  \begin{itemize}
  \tightlist
  \item
    core authoring helper for app-defined primitives
  \end{itemize}
\item
  Graceful lowering:

  \begin{itemize}
  \tightlist
  \item
    rich/static: use the richest honest rendering for the environment
  \item
    pipe: lower to plain, explicit text without hidden styling assumptions
  \item
    accessible: preserve the same semantic meaning in a reading-order-friendly form
  \end{itemize}
\item
  Related families:

  \begin{itemize}
  \tightlist
  \item
    \texttt{n\allowbreak{}o\allowbreak{}t\allowbreak{}e\allowbreak{}(\allowbreak{})}
  \item
    \texttt{b\allowbreak{}a\allowbreak{}d\allowbreak{}g\allowbreak{}e\allowbreak{}(\allowbreak{})}
  \item
    \texttt{m\allowbreak{}a\allowbreak{}r\allowbreak{}k\allowbreak{}d\allowbreak{}o\allowbreak{}w\allowbreak{}n\allowbreak{}(\allowbreak{})}
  \end{itemize}
\item
  Carbon analogue:

  \begin{itemize}
  \tightlist
  \item
    closest to an internal component authoring pattern rather than a shipped Carbon family
  \end{itemize}
\end{itemize}

\section{Data and browsing families}\label{record-019-data-and-browsing-families}

\subsection{Dense comparison}\label{record-019-dense-comparison}

\begin{itemize}
\tightlist
\item
  Family:

  \begin{itemize}
  \tightlist
  \item
    \texttt{t\allowbreak{}a\allowbreak{}b\allowbreak{}l\allowbreak{}e\allowbreak{}(\allowbreak{})}
  \item
    \texttt{t\allowbreak{}a\allowbreak{}b\allowbreak{}l\allowbreak{}e\allowbreak{}S\allowbreak{}u\allowbreak{}r\allowbreak{}f\allowbreak{}a\allowbreak{}c\allowbreak{}e\allowbreak{}(\allowbreak{})}
  \item
    \texttt{n\allowbreak{}a\allowbreak{}v\allowbreak{}i\allowbreak{}g\allowbreak{}a\allowbreak{}b\allowbreak{}l\allowbreak{}e\allowbreak{}T\allowbreak{}a\allowbreak{}b\allowbreak{}l\allowbreak{}e\allowbreak{}(\allowbreak{})}
  \item
    \texttt{n\allowbreak{}a\allowbreak{}v\allowbreak{}i\allowbreak{}g\allowbreak{}a\allowbreak{}b\allowbreak{}l\allowbreak{}e\allowbreak{}T\allowbreak{}a\allowbreak{}b\allowbreak{}l\allowbreak{}e\allowbreak{}S\allowbreak{}u\allowbreak{}r\allowbreak{}f\allowbreak{}a\allowbreak{}c\allowbreak{}e\allowbreak{}(\allowbreak{})}
  \end{itemize}
\item
  Variants:

  \begin{itemize}
  \tightlist
  \item
    \texttt{t\allowbreak{}a\allowbreak{}b\allowbreak{}l\allowbreak{}e\allowbreak{}(\allowbreak{})} for passive comparison in core string output
  \item
    \texttt{t\allowbreak{}a\allowbreak{}b\allowbreak{}l\allowbreak{}e\allowbreak{}(\allowbreak{}\{\allowbreak{} \allowbreak{}v\allowbreak{}a\allowbreak{}r\allowbreak{}i\allowbreak{}a\allowbreak{}n\allowbreak{}t\allowbreak{}:\allowbreak{} \allowbreak{}"\allowbreak{}b\allowbreak{}o\allowbreak{}x\allowbreak{}"\allowbreak{} \allowbreak{}\}\allowbreak{})} for the default Unicode boxed table
  \item
    \texttt{t\allowbreak{}a\allowbreak{}b\allowbreak{}l\allowbreak{}e\allowbreak{}(\allowbreak{}\{\allowbreak{} \allowbreak{}v\allowbreak{}a\allowbreak{}r\allowbreak{}i\allowbreak{}a\allowbreak{}n\allowbreak{}t\allowbreak{}:\allowbreak{} \allowbreak{}"\allowbreak{}r\allowbreak{}u\allowbreak{}l\allowbreak{}e\allowbreak{}d\allowbreak{}"\allowbreak{} \allowbreak{}\}\allowbreak{})} for borderless report output with a heavy header rule and light row separators
  \item
    \texttt{t\allowbreak{}a\allowbreak{}b\allowbreak{}l\allowbreak{}e\allowbreak{}(\allowbreak{}\{\allowbreak{} \allowbreak{}v\allowbreak{}a\allowbreak{}r\allowbreak{}i\allowbreak{}a\allowbreak{}n\allowbreak{}t\allowbreak{}:\allowbreak{} \allowbreak{}"\allowbreak{}h\allowbreak{}e\allowbreak{}a\allowbreak{}d\allowbreak{}e\allowbreak{}r\allowbreak{}-\allowbreak{}r\allowbreak{}u\allowbreak{}l\allowbreak{}e\allowbreak{}"\allowbreak{} \allowbreak{}\}\allowbreak{})} for compact reports with only a header separator
  \item
    \texttt{t\allowbreak{}a\allowbreak{}b\allowbreak{}l\allowbreak{}e\allowbreak{}(\allowbreak{}\{\allowbreak{} \allowbreak{}v\allowbreak{}a\allowbreak{}r\allowbreak{}i\allowbreak{}a\allowbreak{}n\allowbreak{}t\allowbreak{}:\allowbreak{} \allowbreak{}"\allowbreak{}p\allowbreak{}l\allowbreak{}a\allowbreak{}i\allowbreak{}n\allowbreak{}"\allowbreak{} \allowbreak{}\}\allowbreak{})} for quiet space-aligned columns
  \item
    \texttt{t\allowbreak{}a\allowbreak{}b\allowbreak{}l\allowbreak{}e\allowbreak{}(\allowbreak{}\{\allowbreak{} \allowbreak{}v\allowbreak{}a\allowbreak{}r\allowbreak{}i\allowbreak{}a\allowbreak{}n\allowbreak{}t\allowbreak{}:\allowbreak{} \allowbreak{}"\allowbreak{}a\allowbreak{}s\allowbreak{}c\allowbreak{}i\allowbreak{}i\allowbreak{}-\allowbreak{}g\allowbreak{}r\allowbreak{}i\allowbreak{}d\allowbreak{}"\allowbreak{} \allowbreak{}\}\allowbreak{})} for portable \texttt{+\allowbreak{}-\allowbreak{}-\allowbreak{}-\allowbreak{}+} output
  \item
    \texttt{t\allowbreak{}a\allowbreak{}b\allowbreak{}l\allowbreak{}e\allowbreak{}(\allowbreak{}\{\allowbreak{} \allowbreak{}v\allowbreak{}a\allowbreak{}r\allowbreak{}i\allowbreak{}a\allowbreak{}n\allowbreak{}t\allowbreak{}:\allowbreak{} \allowbreak{}"\allowbreak{}m\allowbreak{}a\allowbreak{}r\allowbreak{}k\allowbreak{}d\allowbreak{}o\allowbreak{}w\allowbreak{}n\allowbreak{}"\allowbreak{} \allowbreak{}\}\allowbreak{})} for GitHub-flavored Markdown table text
  \item
    \texttt{t\allowbreak{}a\allowbreak{}b\allowbreak{}l\allowbreak{}e\allowbreak{}(\allowbreak{}\{\allowbreak{} \allowbreak{}v\allowbreak{}a\allowbreak{}r\allowbreak{}i\allowbreak{}a\allowbreak{}n\allowbreak{}t\allowbreak{}:\allowbreak{} \allowbreak{}"\allowbreak{}d\allowbreak{}e\allowbreak{}f\allowbreak{}i\allowbreak{}n\allowbreak{}i\allowbreak{}t\allowbreak{}i\allowbreak{}o\allowbreak{}n\allowbreak{}"\allowbreak{} \allowbreak{}\}\allowbreak{})} for field/value reports
  \item
    \texttt{t\allowbreak{}a\allowbreak{}b\allowbreak{}l\allowbreak{}e\allowbreak{}(\allowbreak{}\{\allowbreak{} \allowbreak{}v\allowbreak{}a\allowbreak{}r\allowbreak{}i\allowbreak{}a\allowbreak{}n\allowbreak{}t\allowbreak{}:\allowbreak{} \allowbreak{}"\allowbreak{}e\allowbreak{}x\allowbreak{}p\allowbreak{}a\allowbreak{}n\allowbreak{}d\allowbreak{}e\allowbreak{}d\allowbreak{}"\allowbreak{} \allowbreak{}\}\allowbreak{})} for psql-style record inspection
  \item
    \texttt{t\allowbreak{}a\allowbreak{}b\allowbreak{}l\allowbreak{}e\allowbreak{}S\allowbreak{}u\allowbreak{}r\allowbreak{}f\allowbreak{}a\allowbreak{}c\allowbreak{}e\allowbreak{}(\allowbreak{})} for passive comparison in V3 surface-first output
  \item
    \texttt{n\allowbreak{}a\allowbreak{}v\allowbreak{}i\allowbreak{}g\allowbreak{}a\allowbreak{}b\allowbreak{}l\allowbreak{}e\allowbreak{}T\allowbreak{}a\allowbreak{}b\allowbreak{}l\allowbreak{}e\allowbreak{}(\allowbreak{})} for keyboard-owned inspection in the TUI layer
  \item
    \texttt{n\allowbreak{}a\allowbreak{}v\allowbreak{}i\allowbreak{}g\allowbreak{}a\allowbreak{}b\allowbreak{}l\allowbreak{}e\allowbreak{}T\allowbreak{}a\allowbreak{}b\allowbreak{}l\allowbreak{}e\allowbreak{}S\allowbreak{}u\allowbreak{}r\allowbreak{}f\allowbreak{}a\allowbreak{}c\allowbreak{}e\allowbreak{}(\allowbreak{})} for keyboard-owned inspection on the structured surface path
  \end{itemize}
\item
  Use when:

  \begin{itemize}
  \tightlist
  \item
    row/column comparison is the main task
  \end{itemize}
\item
  Avoid when:

  \begin{itemize}
  \tightlist
  \item
    hierarchy or dependency structure dominates
  \item
    items are primarily one-dimensional and should read as a list instead
  \end{itemize}
\item
  Content guidance:

  \begin{itemize}
  \tightlist
  \item
    column headers should describe comparable attributes, not narrative explanations
  \item
    cell text should stay compact enough that wrapping does not obscure the comparison task
  \item
    \texttt{t\allowbreak{}a\allowbreak{}b\allowbreak{}l\allowbreak{}e\allowbreak{}(\allowbreak{})} fits human output to the available terminal width by default; use \texttt{l\allowbreak{}a\allowbreak{}y\allowbreak{}o\allowbreak{}u\allowbreak{}t\allowbreak{}:\allowbreak{} \allowbreak{}"\allowbreak{}i\allowbreak{}n\allowbreak{}t\allowbreak{}r\allowbreak{}i\allowbreak{}n\allowbreak{}s\allowbreak{}i\allowbreak{}c\allowbreak{}"} when content-sized output is explicitly needed
  \item
    keep pipe/data serialization explicit with \texttt{p\allowbreak{}i\allowbreak{}p\allowbreak{}e\allowbreak{}F\allowbreak{}o\allowbreak{}r\allowbreak{}m\allowbreak{}a\allowbreak{}t}; TSV remains the default pipe lowering, while CSV, Markdown, and ASCII grid are opt-in
  \item
    use \texttt{n\allowbreak{}a\allowbreak{}v\allowbreak{}i\allowbreak{}g\allowbreak{}a\allowbreak{}b\allowbreak{}l\allowbreak{}e\allowbreak{}T\allowbreak{}a\allowbreak{}b\allowbreak{}l\allowbreak{}e\allowbreak{}S\allowbreak{}u\allowbreak{}r\allowbreak{}f\allowbreak{}a\allowbreak{}c\allowbreak{}e\allowbreak{}(\allowbreak{})} when the table lives inside a rich TUI surface and the row-aware inspection model should remain structured
  \item
    unlike list/picker/palette families, navigable tables should keep row-aware scrolling semantics rather than generic line clipping whenever wrapped rows materially affect comparison
  \item
    if rows wrap so heavily that column comparison is no longer readable, the table should usually collapse to a summary row plus a focused drill-down region instead of pretending dense comparison still works
  \item
    choose \texttt{n\allowbreak{}a\allowbreak{}v\allowbreak{}i\allowbreak{}g\allowbreak{}a\allowbreak{}b\allowbreak{}l\allowbreak{}e\allowbreak{}T\allowbreak{}a\allowbreak{}b\allowbreak{}l\allowbreak{}e\allowbreak{}(\allowbreak{})} / \texttt{n\allowbreak{}a\allowbreak{}v\allowbreak{}i\allowbreak{}g\allowbreak{}a\allowbreak{}b\allowbreak{}l\allowbreak{}e\allowbreak{}T\allowbreak{}a\allowbreak{}b\allowbreak{}l\allowbreak{}e\allowbreak{}S\allowbreak{}u\allowbreak{}r\allowbreak{}f\allowbreak{}a\allowbreak{}c\allowbreak{}e\allowbreak{}(\allowbreak{})} only when keyboard traversal is actually part of the job; otherwise keep the table passive
  \end{itemize}
\item
  Ownership:

  \begin{itemize}
  \tightlist
  \item
    core plus TUI interaction layer
  \end{itemize}
\item
  Graceful lowering:

  \begin{itemize}
  \tightlist
  \item
    rich: passive table or focused navigable inspection depending on variant
  \item
    static: retain tabular comparison when width allows, otherwise favor honest wrapped rows
  \item
    pipe: lower to textual row/column output without hidden clipping
  \item
    accessible: preserve headers, row labels, and comparison semantics explicitly
  \end{itemize}
\item
  Related families:

  \begin{itemize}
  \tightlist
  \item
    \texttt{b\allowbreak{}r\allowbreak{}o\allowbreak{}w\allowbreak{}s\allowbreak{}a\allowbreak{}b\allowbreak{}l\allowbreak{}e\allowbreak{}L\allowbreak{}i\allowbreak{}s\allowbreak{}t\allowbreak{}(\allowbreak{})}
  \item
    \texttt{t\allowbreak{}r\allowbreak{}e\allowbreak{}e\allowbreak{}(\allowbreak{})}
  \item
    \texttt{n\allowbreak{}a\allowbreak{}v\allowbreak{}i\allowbreak{}g\allowbreak{}a\allowbreak{}b\allowbreak{}l\allowbreak{}e\allowbreak{}T\allowbreak{}a\allowbreak{}b\allowbreak{}l\allowbreak{}e\allowbreak{}(\allowbreak{})}
  \end{itemize}
\item
  Carbon analogue:

  \begin{itemize}
  \tightlist
  \item
    data table
  \end{itemize}
\end{itemize}

\subsection{Hierarchy}\label{record-019-hierarchy}

\begin{itemize}
\tightlist
\item
  Family:

  \begin{itemize}
  \tightlist
  \item
    \texttt{t\allowbreak{}r\allowbreak{}e\allowbreak{}e\allowbreak{}(\allowbreak{})}
  \item
    \texttt{f\allowbreak{}i\allowbreak{}l\allowbreak{}e\allowbreak{}P\allowbreak{}i\allowbreak{}c\allowbreak{}k\allowbreak{}e\allowbreak{}r\allowbreak{}(\allowbreak{})}
  \item
    \texttt{f\allowbreak{}i\allowbreak{}l\allowbreak{}e\allowbreak{}P\allowbreak{}i\allowbreak{}c\allowbreak{}k\allowbreak{}e\allowbreak{}r\allowbreak{}S\allowbreak{}u\allowbreak{}r\allowbreak{}f\allowbreak{}a\allowbreak{}c\allowbreak{}e\allowbreak{}(\allowbreak{})}
  \end{itemize}
\item
  Variants:

  \begin{itemize}
  \tightlist
  \item
    static hierarchy
  \item
    interactive filesystem browser
  \end{itemize}
\item
  Use when:

  \begin{itemize}
  \tightlist
  \item
    parent/child nesting is the mental model
  \end{itemize}
\item
  Avoid when:

  \begin{itemize}
  \tightlist
  \item
    multiple parents or graph dependencies matter
  \item
    the interaction is really filesystem traversal, where \texttt{f\allowbreak{}i\allowbreak{}l\allowbreak{}e\allowbreak{}P\allowbreak{}i\allowbreak{}c\allowbreak{}k\allowbreak{}e\allowbreak{}r\allowbreak{}(\allowbreak{})} / \texttt{f\allowbreak{}i\allowbreak{}l\allowbreak{}e\allowbreak{}P\allowbreak{}i\allowbreak{}c\allowbreak{}k\allowbreak{}e\allowbreak{}r\allowbreak{}S\allowbreak{}u\allowbreak{}r\allowbreak{}f\allowbreak{}a\allowbreak{}c\allowbreak{}e\allowbreak{}(\allowbreak{})} are more honest
  \item
    dense comparison matters more than nesting
  \end{itemize}
\item
  Content guidance:

  \begin{itemize}
  \tightlist
  \item
    labels should make nesting meaningful even when indentation is flattened or linearized
  \item
    path context should be explicit when many siblings have similar names
  \item
    use \texttt{f\allowbreak{}i\allowbreak{}l\allowbreak{}e\allowbreak{}P\allowbreak{}i\allowbreak{}c\allowbreak{}k\allowbreak{}e\allowbreak{}r\allowbreak{}S\allowbreak{}u\allowbreak{}r\allowbreak{}f\allowbreak{}a\allowbreak{}c\allowbreak{}e\allowbreak{}(\allowbreak{})} when the picker lives inside a rich TUI pane and should inherit shared viewport masking semantics; keep \texttt{f\allowbreak{}i\allowbreak{}l\allowbreak{}e\allowbreak{}P\allowbreak{}i\allowbreak{}c\allowbreak{}k\allowbreak{}e\allowbreak{}r\allowbreak{}(\allowbreak{})} as the explicit text-lowering path
  \item
    directories should remain visually distinct from files, and moving to a parent or child path should be explicit rather than hidden behind overloaded actions
  \item
    \texttt{t\allowbreak{}r\allowbreak{}e\allowbreak{}e\allowbreak{}(\allowbreak{})} should stay focused on passive hierarchy display; if arbitrary hierarchy needs keyboard-owned non-filesystem inspection later, that should become an explicit interactive tree family instead of overloading \texttt{t\allowbreak{}r\allowbreak{}e\allowbreak{}e\allowbreak{}(\allowbreak{})} itself
  \end{itemize}
\item
  Ownership:

  \begin{itemize}
  \tightlist
  \item
    core plus TUI interaction layer
  \end{itemize}
\item
  Graceful lowering:

  \begin{itemize}
  \tightlist
  \item
    rich: render hierarchy with optional interactive file/navigation layer
  \item
    static: keep visible nesting and path structure honestly
  \item
    pipe: lower naturally to textual indented hierarchy/path lists
  \item
    accessible: preserve parent/child relationships in linear text
  \end{itemize}
\item
  Related families:

  \begin{itemize}
  \tightlist
  \item
    \texttt{b\allowbreak{}r\allowbreak{}o\allowbreak{}w\allowbreak{}s\allowbreak{}a\allowbreak{}b\allowbreak{}l\allowbreak{}e\allowbreak{}L\allowbreak{}i\allowbreak{}s\allowbreak{}t\allowbreak{}(\allowbreak{})}
  \item
    \texttt{d\allowbreak{}a\allowbreak{}g\allowbreak{}(\allowbreak{})}
  \item
    \texttt{f\allowbreak{}i\allowbreak{}l\allowbreak{}e\allowbreak{}P\allowbreak{}i\allowbreak{}c\allowbreak{}k\allowbreak{}e\allowbreak{}r\allowbreak{}(\allowbreak{})}
  \end{itemize}
\item
  Carbon analogue:

  \begin{itemize}
  \tightlist
  \item
    tree view / file browser pattern
  \end{itemize}
\end{itemize}

\subsection{Lists for exploration}\label{record-019-lists-for-exploration}

\begin{itemize}
\tightlist
\item
  Family:

  \begin{itemize}
  \tightlist
  \item
    \texttt{e\allowbreak{}n\allowbreak{}u\allowbreak{}m\allowbreak{}e\allowbreak{}r\allowbreak{}a\allowbreak{}t\allowbreak{}e\allowbreak{}d\allowbreak{}L\allowbreak{}i\allowbreak{}s\allowbreak{}t\allowbreak{}(\allowbreak{})}
  \item
    \texttt{b\allowbreak{}r\allowbreak{}o\allowbreak{}w\allowbreak{}s\allowbreak{}a\allowbreak{}b\allowbreak{}l\allowbreak{}e\allowbreak{}L\allowbreak{}i\allowbreak{}s\allowbreak{}t\allowbreak{}(\allowbreak{})}
  \item
    \texttt{b\allowbreak{}r\allowbreak{}o\allowbreak{}w\allowbreak{}s\allowbreak{}a\allowbreak{}b\allowbreak{}l\allowbreak{}e\allowbreak{}L\allowbreak{}i\allowbreak{}s\allowbreak{}t\allowbreak{}S\allowbreak{}u\allowbreak{}r\allowbreak{}f\allowbreak{}a\allowbreak{}c\allowbreak{}e\allowbreak{}(\allowbreak{})}
  \end{itemize}
\item
  Variants:

  \begin{itemize}
  \tightlist
  \item
    passive list
  \item
    interactive browsable list
  \end{itemize}
\item
  Use when:

  \begin{itemize}
  \tightlist
  \item
    one-dimensional items are being scanned or explored
  \end{itemize}
\item
  Avoid when:

  \begin{itemize}
  \tightlist
  \item
    columns or hierarchy carry the meaning
  \end{itemize}
\item
  Content guidance:

  \begin{itemize}
  \tightlist
  \item
    list rows should begin with the most distinguishing label or identifier
  \item
    descriptions should help the scan task, not turn each row into a paragraph
  \item
    use \texttt{b\allowbreak{}r\allowbreak{}o\allowbreak{}w\allowbreak{}s\allowbreak{}a\allowbreak{}b\allowbreak{}l\allowbreak{}e\allowbreak{}L\allowbreak{}i\allowbreak{}s\allowbreak{}t\allowbreak{}S\allowbreak{}u\allowbreak{}r\allowbreak{}f\allowbreak{}a\allowbreak{}c\allowbreak{}e\allowbreak{}(\allowbreak{})} when the list belongs inside a rich TUI region and should share viewport masking semantics with pagers and focus panes; keep \texttt{b\allowbreak{}r\allowbreak{}o\allowbreak{}w\allowbreak{}s\allowbreak{}a\allowbreak{}b\allowbreak{}l\allowbreak{}e\allowbreak{}L\allowbreak{}i\allowbreak{}s\allowbreak{}t\allowbreak{}(\allowbreak{})} as the explicit text-lowering path
  \item
    if the next user action is ``run a command'' rather than ``inspect a record,'' move up to \texttt{c\allowbreak{}o\allowbreak{}m\allowbreak{}m\allowbreak{}a\allowbreak{}n\allowbreak{}d\allowbreak{}P\allowbreak{}a\allowbreak{}l\allowbreak{}e\allowbreak{}t\allowbreak{}t\allowbreak{}e\allowbreak{}(\allowbreak{})} instead of making the list pretend to be an action launcher
  \end{itemize}
\item
  Ownership:

  \begin{itemize}
  \tightlist
  \item
    core plus TUI interaction layer
  \end{itemize}
\item
  Graceful lowering:

  \begin{itemize}
  \tightlist
  \item
    rich: passive or keyboard-browsable list with optional descriptions
  \item
    static: keep readable ordered list snapshot
  \item
    pipe: lower naturally to plain item list text
  \item
    accessible: preserve item ordering, active selection, and descriptions explicitly
  \end{itemize}
\item
  Related families:

  \begin{itemize}
  \tightlist
  \item
    \texttt{t\allowbreak{}a\allowbreak{}b\allowbreak{}l\allowbreak{}e\allowbreak{}(\allowbreak{})}
  \item
    \texttt{t\allowbreak{}r\allowbreak{}e\allowbreak{}e\allowbreak{}(\allowbreak{})}
  \item
    \texttt{c\allowbreak{}o\allowbreak{}m\allowbreak{}m\allowbreak{}a\allowbreak{}n\allowbreak{}d\allowbreak{}P\allowbreak{}a\allowbreak{}l\allowbreak{}e\allowbreak{}t\allowbreak{}t\allowbreak{}e\allowbreak{}(\allowbreak{})}
  \end{itemize}
\item
  Carbon analogue:

  \begin{itemize}
  \tightlist
  \item
    ordered/unordered list / selectable list
  \end{itemize}
\end{itemize}

\subsection{Temporal or dependency views}\label{record-019-temporal-or-dependency-views}

\begin{itemize}
\tightlist
\item
  Family:

  \begin{itemize}
  \tightlist
  \item
    \texttt{t\allowbreak{}i\allowbreak{}m\allowbreak{}e\allowbreak{}l\allowbreak{}i\allowbreak{}n\allowbreak{}e\allowbreak{}(\allowbreak{})}
  \item
    \texttt{d\allowbreak{}a\allowbreak{}g\allowbreak{}(\allowbreak{})}
  \item
    \texttt{d\allowbreak{}a\allowbreak{}g\allowbreak{}S\allowbreak{}l\allowbreak{}i\allowbreak{}c\allowbreak{}e\allowbreak{}(\allowbreak{})}
  \item
    \texttt{d\allowbreak{}a\allowbreak{}g\allowbreak{}S\allowbreak{}t\allowbreak{}a\allowbreak{}t\allowbreak{}s\allowbreak{}(\allowbreak{})}
  \item
    \texttt{d\allowbreak{}a\allowbreak{}g\allowbreak{}P\allowbreak{}a\allowbreak{}n\allowbreak{}e\allowbreak{}(\allowbreak{})}
  \end{itemize}
\item
  Variants:

  \begin{itemize}
  \tightlist
  \item
    chronological
  \item
    dependency graph
  \item
    focused fragment
  \item
    metrics
  \item
    interactive graph inspection
  \end{itemize}
\item
  Use when:

  \begin{itemize}
  \tightlist
  \item
    time or dependency is the actual structure
  \item
    readers need to answer ``what happened next?'' or ``what depends on this?''
  \end{itemize}
\item
  Avoid when:

  \begin{itemize}
  \tightlist
  \item
    a plain table or tree is enough
  \item
    the graph is mostly decorative architecture wallpaper
  \item
    a summary metric or local fragment would answer the question more honestly than the whole structure
  \end{itemize}
\item
  Content guidance:

  \begin{itemize}
  \tightlist
  \item
    events and nodes should expose the causal or temporal relationship, not just labels
  \item
    annotations should stay lightweight enough that the structure remains legible
  \item
    use \texttt{t\allowbreak{}i\allowbreak{}m\allowbreak{}e\allowbreak{}l\allowbreak{}i\allowbreak{}n\allowbreak{}e\allowbreak{}(\allowbreak{})} when chronology is the actual reading path, not as a decorative status list
  \item
    if timeline rows become so dense that the user is comparing attributes instead of following sequence, collapse to summary milestones plus drill-down or move to a table/log view
  \item
    use \texttt{d\allowbreak{}a\allowbreak{}g\allowbreak{}(\allowbreak{})} / \texttt{d\allowbreak{}a\allowbreak{}g\allowbreak{}P\allowbreak{}a\allowbreak{}n\allowbreak{}e\allowbreak{}(\allowbreak{})} when causality, dependency, or multi-parent structure matters more than order
  \item
    use \texttt{d\allowbreak{}a\allowbreak{}g\allowbreak{}P\allowbreak{}a\allowbreak{}n\allowbreak{}e\allowbreak{}(\allowbreak{})} when graph inspection is an active task and keyboard-owned navigation, focus, and scrolling are part of the job; keep plain \texttt{d\allowbreak{}a\allowbreak{}g\allowbreak{}(\allowbreak{})} for passive explanation
  \item
    use \texttt{d\allowbreak{}a\allowbreak{}g\allowbreak{}S\allowbreak{}l\allowbreak{}i\allowbreak{}c\allowbreak{}e\allowbreak{}(\allowbreak{})} when the honest scope is a neighborhood, ancestor chain, or descendant chain and the full graph would exceed the reader's cognitive budget
  \item
    use \texttt{d\allowbreak{}a\allowbreak{}g\allowbreak{}S\allowbreak{}t\allowbreak{}a\allowbreak{}t\allowbreak{}s\allowbreak{}(\allowbreak{})} when graph health, breadth, depth, or duplication risk matters more than visual shape
  \item
    if most nodes have only one parent and the reader still consumes them in sequence, \texttt{t\allowbreak{}i\allowbreak{}m\allowbreak{}e\allowbreak{}l\allowbreak{}i\allowbreak{}n\allowbreak{}e\allowbreak{}(\allowbreak{})} or \texttt{t\allowbreak{}r\allowbreak{}e\allowbreak{}e\allowbreak{}(\allowbreak{})} is usually clearer than a DAG
  \item
    audit-trail timelines should preserve event ordering and consequence clearly, but avoid stuffing each row with enough metadata that the temporal story disappears
  \end{itemize}
\item
  Ownership:

  \begin{itemize}
  \tightlist
  \item
    core plus TUI interaction layer
  \end{itemize}
\item
  Graceful lowering:

  \begin{itemize}
  \tightlist
  \item
    rich: preserve chronology or graph shape with optional focused inspection
  \item
    static: keep time/dependency order honestly even if the spatial rendering simplifies
  \item
    pipe: lower to ordered event lists or textual dependency traces
  \item
    accessible: linearize the structure while preserving temporal/dependency meaning explicitly
  \end{itemize}
\item
  Related families:

  \begin{itemize}
  \tightlist
  \item
    \texttt{t\allowbreak{}a\allowbreak{}b\allowbreak{}l\allowbreak{}e\allowbreak{}(\allowbreak{})}
  \item
    \texttt{t\allowbreak{}r\allowbreak{}e\allowbreak{}e\allowbreak{}(\allowbreak{})}
  \item
    \texttt{l\allowbreak{}o\allowbreak{}g\allowbreak{}(\allowbreak{})}
  \item
    \texttt{d\allowbreak{}a\allowbreak{}g\allowbreak{}P\allowbreak{}a\allowbreak{}n\allowbreak{}e\allowbreak{}(\allowbreak{})}
  \end{itemize}
\item
  Carbon analogue:

  \begin{itemize}
  \tightlist
  \item
    timeline has a close analogue; DAG does not and should be treated as a specialized Bijou family
  \end{itemize}
\end{itemize}

\subsection{Viewport masking and scrollable inspection panes}\label{record-019-viewport-masking-and-scrollable-inspection-panes}

\begin{itemize}
\tightlist
\item
  Family:

  \begin{itemize}
  \tightlist
  \item
    \texttt{v\allowbreak{}i\allowbreak{}e\allowbreak{}w\allowbreak{}p\allowbreak{}o\allowbreak{}r\allowbreak{}t\allowbreak{}(\allowbreak{})}
  \item
    \texttt{v\allowbreak{}i\allowbreak{}e\allowbreak{}w\allowbreak{}p\allowbreak{}o\allowbreak{}r\allowbreak{}t\allowbreak{}S\allowbreak{}u\allowbreak{}r\allowbreak{}f\allowbreak{}a\allowbreak{}c\allowbreak{}e\allowbreak{}(\allowbreak{})}
  \item
    \texttt{p\allowbreak{}a\allowbreak{}g\allowbreak{}e\allowbreak{}r\allowbreak{}(\allowbreak{})}
  \item
    \texttt{p\allowbreak{}a\allowbreak{}g\allowbreak{}e\allowbreak{}r\allowbreak{}S\allowbreak{}u\allowbreak{}r\allowbreak{}f\allowbreak{}a\allowbreak{}c\allowbreak{}e\allowbreak{}(\allowbreak{})}
  \item
    \texttt{f\allowbreak{}o\allowbreak{}c\allowbreak{}u\allowbreak{}s\allowbreak{}A\allowbreak{}r\allowbreak{}e\allowbreak{}a\allowbreak{}(\allowbreak{})}
  \item
    \texttt{f\allowbreak{}o\allowbreak{}c\allowbreak{}u\allowbreak{}s\allowbreak{}A\allowbreak{}r\allowbreak{}e\allowbreak{}a\allowbreak{}S\allowbreak{}u\allowbreak{}r\allowbreak{}f\allowbreak{}a\allowbreak{}c\allowbreak{}e\allowbreak{}(\allowbreak{})}
  \end{itemize}
\item
  Variants:

  \begin{itemize}
  \tightlist
  \item
    pure scroll mask
  \item
    linear pager
  \item
    focused pane with gutter
  \item
    text-lowering path
  \item
    surface-native path
  \end{itemize}
\item
  Use when:

  \begin{itemize}
  \tightlist
  \item
    a bounded pane is for reading, reviewing, or scrolling through content rather than navigating a richer table, tree, or graph structure
  \end{itemize}
\item
  Avoid when:

  \begin{itemize}
  \tightlist
  \item
    the content needs row/column comparison, hierarchy, or domain-specific navigation semantics
  \end{itemize}
\item
  Ownership:

  \begin{itemize}
  \tightlist
  \item
    TUI
  \end{itemize}
\item
  Content guidance:

  \begin{itemize}
  \tightlist
  \item
    use \texttt{v\allowbreak{}i\allowbreak{}e\allowbreak{}w\allowbreak{}p\allowbreak{}o\allowbreak{}r\allowbreak{}t\allowbreak{}S\allowbreak{}u\allowbreak{}r\allowbreak{}f\allowbreak{}a\allowbreak{}c\allowbreak{}e\allowbreak{}(\allowbreak{})} as the base masking primitive when any existing \texttt{S\allowbreak{}u\allowbreak{}r\allowbreak{}f\allowbreak{}a\allowbreak{}c\allowbreak{}e} or layout-backed region simply needs bounded overflow scrolling
  \item
    keep \texttt{v\allowbreak{}i\allowbreak{}e\allowbreak{}w\allowbreak{}p\allowbreak{}o\allowbreak{}r\allowbreak{}t\allowbreak{}(\allowbreak{})} as the explicit text-lowering path for string-first output, not the default teaching path for rich TUI composition
  \item
    use \texttt{p\allowbreak{}a\allowbreak{}g\allowbreak{}e\allowbreak{}r\allowbreak{}(\allowbreak{})} / \texttt{p\allowbreak{}a\allowbreak{}g\allowbreak{}e\allowbreak{}r\allowbreak{}S\allowbreak{}u\allowbreak{}r\allowbreak{}f\allowbreak{}a\allowbreak{}c\allowbreak{}e\allowbreak{}(\allowbreak{})} for long linear text where the status line and current line position matter
  \item
    use \texttt{f\allowbreak{}o\allowbreak{}c\allowbreak{}u\allowbreak{}s\allowbreak{}A\allowbreak{}r\allowbreak{}e\allowbreak{}a\allowbreak{}(\allowbreak{})} / \texttt{f\allowbreak{}o\allowbreak{}c\allowbreak{}u\allowbreak{}s\allowbreak{}A\allowbreak{}r\allowbreak{}e\allowbreak{}a\allowbreak{}S\allowbreak{}u\allowbreak{}r\allowbreak{}f\allowbreak{}a\allowbreak{}c\allowbreak{}e\allowbreak{}(\allowbreak{})} when the pane participates in a larger workspace and needs explicit focus ownership
  \item
    prefer the surface-native path when the pane body is already composed from \texttt{S\allowbreak{}u\allowbreak{}r\allowbreak{}f\allowbreak{}a\allowbreak{}c\allowbreak{}e} content; keep the string path as explicit lowering, not the default teaching path
  \item
    gutter chrome should communicate focus and workspace ownership, not carry primary content meaning
  \item
    \texttt{b\allowbreak{}r\allowbreak{}o\allowbreak{}w\allowbreak{}s\allowbreak{}a\allowbreak{}b\allowbreak{}l\allowbreak{}e\allowbreak{}L\allowbreak{}i\allowbreak{}s\allowbreak{}t\allowbreak{}S\allowbreak{}u\allowbreak{}r\allowbreak{}f\allowbreak{}a\allowbreak{}c\allowbreak{}e\allowbreak{}(\allowbreak{})}, \texttt{f\allowbreak{}i\allowbreak{}l\allowbreak{}e\allowbreak{}P\allowbreak{}i\allowbreak{}c\allowbreak{}k\allowbreak{}e\allowbreak{}r\allowbreak{}S\allowbreak{}u\allowbreak{}r\allowbreak{}f\allowbreak{}a\allowbreak{}c\allowbreak{}e\allowbreak{}(\allowbreak{})}, and \texttt{c\allowbreak{}o\allowbreak{}m\allowbreak{}m\allowbreak{}a\allowbreak{}n\allowbreak{}d\allowbreak{}P\allowbreak{}a\allowbreak{}l\allowbreak{}e\allowbreak{}t\allowbreak{}t\allowbreak{}e\allowbreak{}S\allowbreak{}u\allowbreak{}r\allowbreak{}f\allowbreak{}a\allowbreak{}c\allowbreak{}e\allowbreak{}(\allowbreak{})} now converge on viewport semantics; \texttt{n\allowbreak{}a\allowbreak{}v\allowbreak{}i\allowbreak{}g\allowbreak{}a\allowbreak{}b\allowbreak{}l\allowbreak{}e\allowbreak{}T\allowbreak{}a\allowbreak{}b\allowbreak{}l\allowbreak{}e\allowbreak{}(\allowbreak{})} remains the row-aware holdout because wrapped-row comparison still needs a more specialized scroll model
  \item
    viewport masking is the wrong abstraction when the user thinks in semantic rows rather than rendered lines; in that case keep row-aware scroll ownership local to the domain-specific component
  \end{itemize}
\item
  Graceful lowering:

  \begin{itemize}
  \tightlist
  \item
    rich: scrollable pane with honest status/focus chrome and bounded viewport behavior
  \item
    static: visible excerpt or current window with enough status context to explain where the user is
  \item
    pipe: lower to sequential text without pretending the hidden region is still interactively present
  \item
    accessible: linearize the pane content while preserving scroll position or focus context explicitly
  \end{itemize}
\item
  Related families:

  \begin{itemize}
  \tightlist
  \item
    \texttt{v\allowbreak{}i\allowbreak{}e\allowbreak{}w\allowbreak{}p\allowbreak{}o\allowbreak{}r\allowbreak{}t\allowbreak{}(\allowbreak{})}
  \item
    \texttt{v\allowbreak{}i\allowbreak{}e\allowbreak{}w\allowbreak{}p\allowbreak{}o\allowbreak{}r\allowbreak{}t\allowbreak{}S\allowbreak{}u\allowbreak{}r\allowbreak{}f\allowbreak{}a\allowbreak{}c\allowbreak{}e\allowbreak{}(\allowbreak{})}
  \item
    \texttt{d\allowbreak{}a\allowbreak{}g\allowbreak{}P\allowbreak{}a\allowbreak{}n\allowbreak{}e\allowbreak{}(\allowbreak{})}
  \item
    \texttt{c\allowbreak{}r\allowbreak{}e\allowbreak{}a\allowbreak{}t\allowbreak{}e\allowbreak{}F\allowbreak{}r\allowbreak{}a\allowbreak{}m\allowbreak{}e\allowbreak{}d\allowbreak{}A\allowbreak{}p\allowbreak{}p\allowbreak{}(\allowbreak{})}
  \item
    \texttt{b\allowbreak{}r\allowbreak{}o\allowbreak{}w\allowbreak{}s\allowbreak{}a\allowbreak{}b\allowbreak{}l\allowbreak{}e\allowbreak{}L\allowbreak{}i\allowbreak{}s\allowbreak{}t\allowbreak{}(\allowbreak{})}
  \end{itemize}
\item
  Carbon analogue:

  \begin{itemize}
  \tightlist
  \item
    closest to a scrollable content region or code viewer; not a direct Carbon family
  \end{itemize}
\end{itemize}

\section{Overlay and shell families}\label{record-019-overlay-and-shell-families}

\subsection{Overlay primitives}\label{record-019-overlay-primitives}

\begin{itemize}
\tightlist
\item
  Family:

  \begin{itemize}
  \tightlist
  \item
    \texttt{t\allowbreak{}o\allowbreak{}o\allowbreak{}l\allowbreak{}t\allowbreak{}i\allowbreak{}p\allowbreak{}(\allowbreak{})}
  \item
    \texttt{d\allowbreak{}r\allowbreak{}a\allowbreak{}w\allowbreak{}e\allowbreak{}r\allowbreak{}(\allowbreak{})}
  \item
    \texttt{m\allowbreak{}o\allowbreak{}d\allowbreak{}a\allowbreak{}l\allowbreak{}(\allowbreak{})}
  \item
    \texttt{t\allowbreak{}o\allowbreak{}a\allowbreak{}s\allowbreak{}t\allowbreak{}(\allowbreak{})}
  \end{itemize}
\item
  Variants:

  \begin{itemize}
  \tightlist
  \item
    explanatory
  \item
    supplemental
  \item
    blocking
  \item
    transient
  \end{itemize}
\item
  Use when:

  \begin{itemize}
  \tightlist
  \item
    the app needs layered interaction surfaces
  \end{itemize}
\item
  Avoid when:

  \begin{itemize}
  \tightlist
  \item
    the content should just live in the document or shell
  \end{itemize}
\item
  Ownership:

  \begin{itemize}
  \tightlist
  \item
    TUI
  \end{itemize}
\item
  Content guidance:

  \begin{itemize}
  \tightlist
  \item
    \texttt{t\allowbreak{}o\allowbreak{}o\allowbreak{}l\allowbreak{}t\allowbreak{}i\allowbreak{}p\allowbreak{}(\allowbreak{})} content should be short, local, and explanatory, not actionable or scrollable
  \item
    \texttt{d\allowbreak{}r\allowbreak{}a\allowbreak{}w\allowbreak{}e\allowbreak{}r\allowbreak{}(\allowbreak{})} content should support supplemental work, inspection, or side-by-side context without stealing the whole task
  \item
    \texttt{m\allowbreak{}o\allowbreak{}d\allowbreak{}a\allowbreak{}l\allowbreak{}(\allowbreak{})} content should justify blocking the user and should end in a clear decision, confirmation, or review step
  \item
    \texttt{t\allowbreak{}o\allowbreak{}a\allowbreak{}s\allowbreak{}t\allowbreak{}(\allowbreak{})} content should be short-lived and self-contained, not a surrogate notification center
  \item
    if the user may need recall, stacking, or routing, stop composing ad hoc \texttt{t\allowbreak{}o\allowbreak{}a\allowbreak{}s\allowbreak{}t\allowbreak{}(\allowbreak{})} overlays and move up to the notification system
  \item
    if the user still needs the main task visible and interactive, prefer \texttt{d\allowbreak{}r\allowbreak{}a\allowbreak{}w\allowbreak{}e\allowbreak{}r\allowbreak{}(\allowbreak{})} over \texttt{m\allowbreak{}o\allowbreak{}d\allowbreak{}a\allowbreak{}l\allowbreak{}(\allowbreak{})}
  \item
    when rich TUI structure matters, overlays should accept or compose structured \texttt{S\allowbreak{}u\allowbreak{}r\allowbreak{}f\allowbreak{}a\allowbreak{}c\allowbreak{}e} content instead of flattening component content into strings first
  \end{itemize}
\item
  Graceful lowering:

  \begin{itemize}
  \tightlist
  \item
    rich: layered overlay surfaces with blocking or non-blocking behavior depending on variant
  \item
    static: lower to visible in-flow or shell-level snapshots rather than hidden z-order
  \item
    pipe: lower to plain text event or prompt surfaces appropriate to interruption level
  \item
    accessible: linearize overlay content with explicit relationship to the blocked or supplemented context
  \end{itemize}
\item
  Related families:

  \begin{itemize}
  \tightlist
  \item
    notification system
  \item
    \texttt{a\allowbreak{}l\allowbreak{}e\allowbreak{}r\allowbreak{}t\allowbreak{}(\allowbreak{})}
  \item
    \texttt{c\allowbreak{}r\allowbreak{}e\allowbreak{}a\allowbreak{}t\allowbreak{}e\allowbreak{}F\allowbreak{}r\allowbreak{}a\allowbreak{}m\allowbreak{}e\allowbreak{}d\allowbreak{}A\allowbreak{}p\allowbreak{}p\allowbreak{}(\allowbreak{})}
  \end{itemize}
\item
  Carbon analogue:

  \begin{itemize}
  \tightlist
  \item
    tooltip / side panel / modal / toast notification
  \end{itemize}
\end{itemize}

\subsection{Notification system}\label{record-019-notification-system}

\begin{itemize}
\tightlist
\item
  Family:

  \begin{itemize}
  \tightlist
  \item
    \texttt{r\allowbreak{}e\allowbreak{}n\allowbreak{}d\allowbreak{}e\allowbreak{}r\allowbreak{}N\allowbreak{}o\allowbreak{}t\allowbreak{}i\allowbreak{}f\allowbreak{}i\allowbreak{}c\allowbreak{}a\allowbreak{}t\allowbreak{}i\allowbreak{}o\allowbreak{}n\allowbreak{}S\allowbreak{}t\allowbreak{}a\allowbreak{}c\allowbreak{}k\allowbreak{}(\allowbreak{})}
  \item
    \texttt{r\allowbreak{}e\allowbreak{}n\allowbreak{}d\allowbreak{}e\allowbreak{}r\allowbreak{}N\allowbreak{}o\allowbreak{}t\allowbreak{}i\allowbreak{}f\allowbreak{}i\allowbreak{}c\allowbreak{}a\allowbreak{}t\allowbreak{}i\allowbreak{}o\allowbreak{}n\allowbreak{}H\allowbreak{}i\allowbreak{}s\allowbreak{}t\allowbreak{}o\allowbreak{}r\allowbreak{}y\allowbreak{}(\allowbreak{})}
  \item
    \texttt{p\allowbreak{}u\allowbreak{}s\allowbreak{}h\allowbreak{}N\allowbreak{}o\allowbreak{}t\allowbreak{}i\allowbreak{}f\allowbreak{}i\allowbreak{}c\allowbreak{}a\allowbreak{}t\allowbreak{}i\allowbreak{}o\allowbreak{}n\allowbreak{}(\allowbreak{})}
  \item
    \texttt{d\allowbreak{}i\allowbreak{}s\allowbreak{}m\allowbreak{}i\allowbreak{}s\allowbreak{}s\allowbreak{}N\allowbreak{}o\allowbreak{}t\allowbreak{}i\allowbreak{}f\allowbreak{}i\allowbreak{}c\allowbreak{}a\allowbreak{}t\allowbreak{}i\allowbreak{}o\allowbreak{}n\allowbreak{}(\allowbreak{})}
  \item
    \texttt{t\allowbreak{}i\allowbreak{}c\allowbreak{}k\allowbreak{}N\allowbreak{}o\allowbreak{}t\allowbreak{}i\allowbreak{}f\allowbreak{}i\allowbreak{}c\allowbreak{}a\allowbreak{}t\allowbreak{}i\allowbreak{}o\allowbreak{}n\allowbreak{}s\allowbreak{}(\allowbreak{})}
  \item
    framed runtime notification routing
  \end{itemize}
\item
  Variants:

  \begin{itemize}
  \tightlist
  \item
    \texttt{A\allowbreak{}C\allowbreak{}T\allowbreak{}I\allowbreak{}O\allowbreak{}N\allowbreak{}A\allowbreak{}B\allowbreak{}L\allowbreak{}E}
  \item
    \texttt{I\allowbreak{}N\allowbreak{}L\allowbreak{}I\allowbreak{}N\allowbreak{}E}
  \item
    \texttt{T\allowbreak{}O\allowbreak{}A\allowbreak{}S\allowbreak{}T}
  \item
    info, success, warning, error tone variants
  \item
    per-notification placement and duration
  \item
    history/archive view
  \item
    frame-routed runtime notifications
  \end{itemize}
\item
  Use when:

  \begin{itemize}
  \tightlist
  \item
    the app owns transient messaging as a system instead of rendering one ad hoc overlay at a time
  \item
    stacking, placement, routing, or history matter
  \item
    warnings and errors should be reviewable after the moment they first appear
  \end{itemize}
\item
  Avoid when:

  \begin{itemize}
  \tightlist
  \item
    the message should remain in normal page flow as part of the document
  \item
    a single local transient overlay is enough
  \item
    the user must stop and decide before continuing
  \end{itemize}
\item
  Ownership:

  \begin{itemize}
  \tightlist
  \item
    TUI
  \end{itemize}
\item
  Content guidance:

  \begin{itemize}
  \tightlist
  \item
    use notifications for events and follow-up prompts, not durable documents or long-form explanation
  \item
    actionable notifications should expose one obvious next step, not a miniature workflow or form
  \item
    if users may need to revisit prior notices, the archive/history surface is part of the feature, not an optional add-on
  \item
    route app-owned warnings and errors here when recall matters, while still keeping stderr/log output honest where appropriate
  \item
    use placement and duration to control interruption level, not to smuggle unrelated importance into every message
  \end{itemize}
\item
  Graceful lowering:

  \begin{itemize}
  \tightlist
  \item
    rich: stacked notifications with actions, placement, history, and explicit dismissal
  \item
    static: visible current notifications or a truthful archived summary without pretending transient timing still exists
  \item
    pipe: lower to sequential event text or explicit warning/error records
  \item
    accessible: linearize current and archived notices with tone, action, and dismissal state made explicit
  \end{itemize}
\item
  Related families:

  \begin{itemize}
  \tightlist
  \item
    \texttt{t\allowbreak{}o\allowbreak{}a\allowbreak{}s\allowbreak{}t\allowbreak{}(\allowbreak{})}
  \item
    \texttt{a\allowbreak{}l\allowbreak{}e\allowbreak{}r\allowbreak{}t\allowbreak{}(\allowbreak{})}
  \item
    \texttt{m\allowbreak{}o\allowbreak{}d\allowbreak{}a\allowbreak{}l\allowbreak{}(\allowbreak{})}
  \item
    \texttt{d\allowbreak{}r\allowbreak{}a\allowbreak{}w\allowbreak{}e\allowbreak{}r\allowbreak{}(\allowbreak{})}
  \item
    \texttt{c\allowbreak{}r\allowbreak{}e\allowbreak{}a\allowbreak{}t\allowbreak{}e\allowbreak{}F\allowbreak{}r\allowbreak{}a\allowbreak{}m\allowbreak{}e\allowbreak{}d\allowbreak{}A\allowbreak{}p\allowbreak{}p\allowbreak{}(\allowbreak{})}
  \end{itemize}
\item
  Carbon analogue:

  \begin{itemize}
  \tightlist
  \item
    notification / toast notification center pattern
  \end{itemize}
\end{itemize}

\subsection{App shell}\label{record-019-app-shell}

\begin{itemize}
\tightlist
\item
  Family:

  \begin{itemize}
  \tightlist
  \item
    \texttt{c\allowbreak{}r\allowbreak{}e\allowbreak{}a\allowbreak{}t\allowbreak{}e\allowbreak{}F\allowbreak{}r\allowbreak{}a\allowbreak{}m\allowbreak{}e\allowbreak{}d\allowbreak{}A\allowbreak{}p\allowbreak{}p\allowbreak{}(\allowbreak{})}
  \item
    \texttt{s\allowbreak{}t\allowbreak{}a\allowbreak{}t\allowbreak{}u\allowbreak{}s\allowbreak{}B\allowbreak{}a\allowbreak{}r\allowbreak{}(\allowbreak{})}
  \item
    \texttt{s\allowbreak{}t\allowbreak{}a\allowbreak{}t\allowbreak{}u\allowbreak{}s\allowbreak{}B\allowbreak{}a\allowbreak{}r\allowbreak{}S\allowbreak{}u\allowbreak{}r\allowbreak{}f\allowbreak{}a\allowbreak{}c\allowbreak{}e\allowbreak{}(\allowbreak{})}
  \item
    \texttt{c\allowbreak{}o\allowbreak{}m\allowbreak{}m\allowbreak{}a\allowbreak{}n\allowbreak{}d\allowbreak{}P\allowbreak{}a\allowbreak{}l\allowbreak{}e\allowbreak{}t\allowbreak{}t\allowbreak{}e\allowbreak{}(\allowbreak{})}
  \item
    \texttt{c\allowbreak{}o\allowbreak{}m\allowbreak{}m\allowbreak{}a\allowbreak{}n\allowbreak{}d\allowbreak{}P\allowbreak{}a\allowbreak{}l\allowbreak{}e\allowbreak{}t\allowbreak{}t\allowbreak{}e\allowbreak{}S\allowbreak{}u\allowbreak{}r\allowbreak{}f\allowbreak{}a\allowbreak{}c\allowbreak{}e\allowbreak{}(\allowbreak{})}
  \end{itemize}
\item
  Variants:

  \begin{itemize}
  \tightlist
  \item
    shell container, status rail, action discovery
  \end{itemize}
\item
  Use when:

  \begin{itemize}
  \tightlist
  \item
    the app has multiple views, overlays, navigation, and shell chrome
  \end{itemize}
\item
  Avoid when:

  \begin{itemize}
  \tightlist
  \item
    the app is really one screen or one prompt
  \end{itemize}
\item
  Ownership:

  \begin{itemize}
  \tightlist
  \item
    TUI
  \end{itemize}
\item
  Content guidance:

  \begin{itemize}
  \tightlist
  \item
    the shell should frame destinations and workspace state, not become a dumping ground for unrelated metadata
  \item
    use \texttt{s\allowbreak{}t\allowbreak{}a\allowbreak{}t\allowbreak{}u\allowbreak{}s\allowbreak{}B\allowbreak{}a\allowbreak{}r\allowbreak{}S\allowbreak{}u\allowbreak{}r\allowbreak{}f\allowbreak{}a\allowbreak{}c\allowbreak{}e\allowbreak{}(\allowbreak{})} when shell chrome already lives on the structured \texttt{S\allowbreak{}u\allowbreak{}r\allowbreak{}f\allowbreak{}a\allowbreak{}c\allowbreak{}e} path; keep \texttt{s\allowbreak{}t\allowbreak{}a\allowbreak{}t\allowbreak{}u\allowbreak{}s\allowbreak{}B\allowbreak{}a\allowbreak{}r\allowbreak{}(\allowbreak{})} for explicit text output or lowering
  \item
    use \texttt{c\allowbreak{}o\allowbreak{}m\allowbreak{}m\allowbreak{}a\allowbreak{}n\allowbreak{}d\allowbreak{}P\allowbreak{}a\allowbreak{}l\allowbreak{}e\allowbreak{}t\allowbreak{}t\allowbreak{}e\allowbreak{}S\allowbreak{}u\allowbreak{}r\allowbreak{}f\allowbreak{}a\allowbreak{}c\allowbreak{}e\allowbreak{}(\allowbreak{})} when the palette lives inside a rich TUI shell or overlay and should share viewport masking semantics with other panes; keep \texttt{c\allowbreak{}o\allowbreak{}m\allowbreak{}m\allowbreak{}a\allowbreak{}n\allowbreak{}d\allowbreak{}P\allowbreak{}a\allowbreak{}l\allowbreak{}e\allowbreak{}t\allowbreak{}t\allowbreak{}e\allowbreak{}(\allowbreak{})} for explicit text output or lowering
  \item
    status lines should carry concise global context, not replace in-page guidance
  \item
    command palette entries should prefer actions, navigation targets, and recent destinations over field-style data entry
  \item
    group palette items by user intent when possible, and avoid turning the palette into a browseable record list
  \item
    tabs should represent peer destinations or work areas, not disguised command buttons
  \item
    use notifications for events and follow-up, not the status rail
  \item
    use help hints for shortcut discovery and shell scope, not for action execution or event messaging
  \end{itemize}
\item
  Graceful lowering:

  \begin{itemize}
  \tightlist
  \item
    rich: full shell chrome with tabs, palette, overlays, notifications, and workspace regions
  \item
    static: retain the active page and essential shell context without pretending background interactivity exists
  \item
    pipe: lower to current page content plus minimal status/context framing
  \item
    accessible: linearize active shell state, navigation context, and current overlays into one readable flow
  \end{itemize}
\item
  Related families:

  \begin{itemize}
  \tightlist
  \item
    \texttt{s\allowbreak{}t\allowbreak{}a\allowbreak{}t\allowbreak{}u\allowbreak{}s\allowbreak{}B\allowbreak{}a\allowbreak{}r\allowbreak{}(\allowbreak{})}
  \item
    \texttt{h\allowbreak{}e\allowbreak{}l\allowbreak{}p\allowbreak{}S\allowbreak{}h\allowbreak{}o\allowbreak{}r\allowbreak{}t\allowbreak{}S\allowbreak{}u\allowbreak{}r\allowbreak{}f\allowbreak{}a\allowbreak{}c\allowbreak{}e\allowbreak{}(\allowbreak{})}
  \item
    \texttt{t\allowbreak{}a\allowbreak{}b\allowbreak{}s\allowbreak{}(\allowbreak{})}
  \item
    \texttt{c\allowbreak{}o\allowbreak{}m\allowbreak{}m\allowbreak{}a\allowbreak{}n\allowbreak{}d\allowbreak{}P\allowbreak{}a\allowbreak{}l\allowbreak{}e\allowbreak{}t\allowbreak{}t\allowbreak{}e\allowbreak{}(\allowbreak{})}
  \item
    notification system
  \end{itemize}
\item
  Carbon analogue:

  \begin{itemize}
  \tightlist
  \item
    shell, header/footer, command palette
  \end{itemize}
\end{itemize}

\subsection{Keybinding help and shell hints}\label{record-019-keybinding-help-and-shell-hints}

\begin{itemize}
\tightlist
\item
  Family:

  \begin{itemize}
  \tightlist
  \item
    \texttt{c\allowbreak{}r\allowbreak{}e\allowbreak{}a\allowbreak{}t\allowbreak{}e\allowbreak{}K\allowbreak{}e\allowbreak{}y\allowbreak{}M\allowbreak{}a\allowbreak{}p\allowbreak{}(\allowbreak{})}
  \item
    \texttt{h\allowbreak{}e\allowbreak{}l\allowbreak{}p\allowbreak{}V\allowbreak{}i\allowbreak{}e\allowbreak{}w\allowbreak{}(\allowbreak{})}
  \item
    \texttt{h\allowbreak{}e\allowbreak{}l\allowbreak{}p\allowbreak{}V\allowbreak{}i\allowbreak{}e\allowbreak{}w\allowbreak{}S\allowbreak{}u\allowbreak{}r\allowbreak{}f\allowbreak{}a\allowbreak{}c\allowbreak{}e\allowbreak{}(\allowbreak{})}
  \item
    \texttt{h\allowbreak{}e\allowbreak{}l\allowbreak{}p\allowbreak{}S\allowbreak{}h\allowbreak{}o\allowbreak{}r\allowbreak{}t\allowbreak{}(\allowbreak{})}
  \item
    \texttt{h\allowbreak{}e\allowbreak{}l\allowbreak{}p\allowbreak{}S\allowbreak{}h\allowbreak{}o\allowbreak{}r\allowbreak{}t\allowbreak{}S\allowbreak{}u\allowbreak{}r\allowbreak{}f\allowbreak{}a\allowbreak{}c\allowbreak{}e\allowbreak{}(\allowbreak{})}
  \item
    \texttt{h\allowbreak{}e\allowbreak{}l\allowbreak{}p\allowbreak{}F\allowbreak{}o\allowbreak{}r\allowbreak{}(\allowbreak{})}
  \item
    \texttt{h\allowbreak{}e\allowbreak{}l\allowbreak{}p\allowbreak{}F\allowbreak{}o\allowbreak{}r\allowbreak{}S\allowbreak{}u\allowbreak{}r\allowbreak{}f\allowbreak{}a\allowbreak{}c\allowbreak{}e\allowbreak{}(\allowbreak{})}
  \end{itemize}
\item
  Variants:

  \begin{itemize}
  \tightlist
  \item
    grouped reference, single-line hint, filtered subset, explicit text output, structured-surface output
  \end{itemize}
\item
  Use when:

  \begin{itemize}
  \tightlist
  \item
    the app is keyboard-owned and the user needs discoverable shortcuts or a grouped command reference
  \end{itemize}
\item
  Avoid when:

  \begin{itemize}
  \tightlist
  \item
    the controls are already obvious from context and the help text would just restate visible labels
  \end{itemize}
\item
  Ownership:

  \begin{itemize}
  \tightlist
  \item
    TUI
  \end{itemize}
\item
  Content guidance:

  \begin{itemize}
  \tightlist
  \item
    use \texttt{h\allowbreak{}e\allowbreak{}l\allowbreak{}p\allowbreak{}S\allowbreak{}h\allowbreak{}o\allowbreak{}r\allowbreak{}t\allowbreak{}S\allowbreak{}u\allowbreak{}r\allowbreak{}f\allowbreak{}a\allowbreak{}c\allowbreak{}e\allowbreak{}(\allowbreak{})} for shell hints that should stay on the structured \texttt{S\allowbreak{}u\allowbreak{}r\allowbreak{}f\allowbreak{}a\allowbreak{}c\allowbreak{}e} path; keep \texttt{h\allowbreak{}e\allowbreak{}l\allowbreak{}p\allowbreak{}S\allowbreak{}h\allowbreak{}o\allowbreak{}r\allowbreak{}t\allowbreak{}(\allowbreak{})} for explicit text output or lowering
  \item
    use \texttt{h\allowbreak{}e\allowbreak{}l\allowbreak{}p\allowbreak{}V\allowbreak{}i\allowbreak{}e\allowbreak{}w\allowbreak{}S\allowbreak{}u\allowbreak{}r\allowbreak{}f\allowbreak{}a\allowbreak{}c\allowbreak{}e\allowbreak{}(\allowbreak{})} when the full keybinding reference is embedded inside a rich TUI surface; keep \texttt{h\allowbreak{}e\allowbreak{}l\allowbreak{}p\allowbreak{}V\allowbreak{}i\allowbreak{}e\allowbreak{}w\allowbreak{}(\allowbreak{})} for plain text and pipe-oriented help
  \item
    group names should describe jobs (\texttt{N\allowbreak{}a\allowbreak{}v\allowbreak{}i\allowbreak{}g\allowbreak{}a\allowbreak{}t\allowbreak{}i\allowbreak{}o\allowbreak{}n}, \texttt{A\allowbreak{}c\allowbreak{}t\allowbreak{}i\allowbreak{}o\allowbreak{}n\allowbreak{}s}, \texttt{S\allowbreak{}e\allowbreak{}l\allowbreak{}e\allowbreak{}c\allowbreak{}t\allowbreak{}i\allowbreak{}o\allowbreak{}n}) rather than raw input mechanics
  \item
    help text should clarify behavior and scope, not repeat page prose or become a substitute command palette
  \item
    use the command palette for discoverable actions and navigation, and keep help focused on shortcut explanation
  \end{itemize}
\item
  Graceful lowering:

  \begin{itemize}
  \tightlist
  \item
    rich: concise shell hints plus an optional grouped help view or modal
  \item
    static: retain the active hint and the relevant grouped reference without pretending background input is live
  \item
    pipe: lower to plain text shortcut summaries and grouped help blocks
  \item
    accessible: linearize help content into readable sections with explicit scope and action labels
  \end{itemize}
\item
  Related families:

  \begin{itemize}
  \tightlist
  \item
    \texttt{c\allowbreak{}r\allowbreak{}e\allowbreak{}a\allowbreak{}t\allowbreak{}e\allowbreak{}F\allowbreak{}r\allowbreak{}a\allowbreak{}m\allowbreak{}e\allowbreak{}d\allowbreak{}A\allowbreak{}p\allowbreak{}p\allowbreak{}(\allowbreak{})}
  \item
    \texttt{s\allowbreak{}t\allowbreak{}a\allowbreak{}t\allowbreak{}u\allowbreak{}s\allowbreak{}B\allowbreak{}a\allowbreak{}r\allowbreak{}S\allowbreak{}u\allowbreak{}r\allowbreak{}f\allowbreak{}a\allowbreak{}c\allowbreak{}e\allowbreak{}(\allowbreak{})}
  \item
    \texttt{c\allowbreak{}o\allowbreak{}m\allowbreak{}m\allowbreak{}a\allowbreak{}n\allowbreak{}d\allowbreak{}P\allowbreak{}a\allowbreak{}l\allowbreak{}e\allowbreak{}t\allowbreak{}t\allowbreak{}e\allowbreak{}(\allowbreak{})}
  \end{itemize}
\item
  Carbon analogue:

  \begin{itemize}
  \tightlist
  \item
    keyboard shortcuts help, command reference, hotkey cheat sheet
  \end{itemize}
\end{itemize}

\subsection{Workspace layout}\label{record-019-workspace-layout}

\begin{itemize}
\tightlist
\item
  Family:

  \begin{itemize}
  \tightlist
  \item
    \texttt{s\allowbreak{}p\allowbreak{}l\allowbreak{}i\allowbreak{}t\allowbreak{}P\allowbreak{}a\allowbreak{}n\allowbreak{}e\allowbreak{}(\allowbreak{})}
  \item
    \texttt{s\allowbreak{}p\allowbreak{}l\allowbreak{}i\allowbreak{}t\allowbreak{}P\allowbreak{}a\allowbreak{}n\allowbreak{}e\allowbreak{}S\allowbreak{}u\allowbreak{}r\allowbreak{}f\allowbreak{}a\allowbreak{}c\allowbreak{}e\allowbreak{}(\allowbreak{})}
  \item
    \texttt{g\allowbreak{}r\allowbreak{}i\allowbreak{}d\allowbreak{}(\allowbreak{})}
  \item
    \texttt{g\allowbreak{}r\allowbreak{}i\allowbreak{}d\allowbreak{}S\allowbreak{}u\allowbreak{}r\allowbreak{}f\allowbreak{}a\allowbreak{}c\allowbreak{}e\allowbreak{}(\allowbreak{})}
  \item
    \texttt{f\allowbreak{}l\allowbreak{}e\allowbreak{}x\allowbreak{}(\allowbreak{})}
  \item
    \texttt{v\allowbreak{}s\allowbreak{}t\allowbreak{}a\allowbreak{}c\allowbreak{}k\allowbreak{}(\allowbreak{})}
  \item
    \texttt{h\allowbreak{}s\allowbreak{}t\allowbreak{}a\allowbreak{}c\allowbreak{}k\allowbreak{}(\allowbreak{})}
  \item
    \texttt{p\allowbreak{}l\allowbreak{}a\allowbreak{}c\allowbreak{}e\allowbreak{}(\allowbreak{})}
  \end{itemize}
\item
  Variants:

  \begin{itemize}
  \tightlist
  \item
    user-resizable, named-area, compositional layout primitives with structured-surface and explicit text-lowering paths
  \end{itemize}
\item
  Use when:

  \begin{itemize}
  \tightlist
  \item
    spatial arrangement materially helps the task
  \end{itemize}
\item
  Avoid when:

  \begin{itemize}
  \tightlist
  \item
    a sequential flow would be simpler and more legible
  \end{itemize}
\item
  Ownership:

  \begin{itemize}
  \tightlist
  \item
    TUI
  \end{itemize}
\item
  Content guidance:

  \begin{itemize}
  \tightlist
  \item
    use \texttt{s\allowbreak{}p\allowbreak{}l\allowbreak{}i\allowbreak{}t\allowbreak{}P\allowbreak{}a\allowbreak{}n\allowbreak{}e\allowbreak{}S\allowbreak{}u\allowbreak{}r\allowbreak{}f\allowbreak{}a\allowbreak{}c\allowbreak{}e\allowbreak{}(\allowbreak{})} when the user benefits from explicit primary-versus-secondary context and the panes are already rich TUI surfaces; keep \texttt{s\allowbreak{}p\allowbreak{}l\allowbreak{}i\allowbreak{}t\allowbreak{}P\allowbreak{}a\allowbreak{}n\allowbreak{}e\allowbreak{}(\allowbreak{})} for explicit text composition or lowering
  \item
    use \texttt{g\allowbreak{}r\allowbreak{}i\allowbreak{}d\allowbreak{}S\allowbreak{}u\allowbreak{}r\allowbreak{}f\allowbreak{}a\allowbreak{}c\allowbreak{}e\allowbreak{}(\allowbreak{})} when multiple stable regions deserve simultaneous visibility and the regions are already rich TUI surfaces; keep \texttt{g\allowbreak{}r\allowbreak{}i\allowbreak{}d\allowbreak{}(\allowbreak{})} for explicit text composition or lowering
  \item
    use \texttt{f\allowbreak{}l\allowbreak{}e\allowbreak{}x\allowbreak{}(\allowbreak{})} / \texttt{v\allowbreak{}s\allowbreak{}t\allowbreak{}a\allowbreak{}c\allowbreak{}k\allowbreak{}S\allowbreak{}u\allowbreak{}r\allowbreak{}f\allowbreak{}a\allowbreak{}c\allowbreak{}e\allowbreak{}(\allowbreak{})} / \texttt{h\allowbreak{}s\allowbreak{}t\allowbreak{}a\allowbreak{}c\allowbreak{}k\allowbreak{}S\allowbreak{}u\allowbreak{}r\allowbreak{}f\allowbreak{}a\allowbreak{}c\allowbreak{}e\allowbreak{}(\allowbreak{})} / \texttt{p\allowbreak{}l\allowbreak{}a\allowbreak{}c\allowbreak{}e\allowbreak{}S\allowbreak{}u\allowbreak{}r\allowbreak{}f\allowbreak{}a\allowbreak{}c\allowbreak{}e\allowbreak{}(\allowbreak{})} to keep rich TUI composition structured; keep \texttt{v\allowbreak{}s\allowbreak{}t\allowbreak{}a\allowbreak{}c\allowbreak{}k\allowbreak{}(\allowbreak{})} / \texttt{h\allowbreak{}s\allowbreak{}t\allowbreak{}a\allowbreak{}c\allowbreak{}k\allowbreak{}(\allowbreak{})} / \texttt{p\allowbreak{}l\allowbreak{}a\allowbreak{}c\allowbreak{}e\allowbreak{}(\allowbreak{})} for explicit text composition or lowering
  \item
    region titles and borders should explain job and hierarchy, not merely expose geometry
  \end{itemize}
\item
  Graceful lowering:

  \begin{itemize}
  \tightlist
  \item
    rich: keep spatial relationships and resizable or placed regions where they materially help
  \item
    static: retain simplified spatial grouping when possible without fake interactivity
  \item
    pipe: lower to sensible sequential content order
  \item
    accessible: linearize regions in a predictable reading order with section labels
  \end{itemize}
\item
  Related families:

  \begin{itemize}
  \tightlist
  \item
    \texttt{b\allowbreak{}o\allowbreak{}x\allowbreak{}(\allowbreak{})}
  \item
    \texttt{c\allowbreak{}r\allowbreak{}e\allowbreak{}a\allowbreak{}t\allowbreak{}e\allowbreak{}F\allowbreak{}r\allowbreak{}a\allowbreak{}m\allowbreak{}e\allowbreak{}d\allowbreak{}A\allowbreak{}p\allowbreak{}p\allowbreak{}(\allowbreak{})}
  \item
    \texttt{f\allowbreak{}o\allowbreak{}c\allowbreak{}u\allowbreak{}s\allowbreak{}A\allowbreak{}r\allowbreak{}e\allowbreak{}a\allowbreak{}(\allowbreak{})}
  \end{itemize}
\item
  Carbon analogue:

  \begin{itemize}
  \tightlist
  \item
    layout grid / split layout / stack primitives
  \end{itemize}
\end{itemize}

\subsection{Data visualization}\label{record-019-data-visualization}

\begin{itemize}
\tightlist
\item
  Family:

  \begin{itemize}
  \tightlist
  \item
    \texttt{s\allowbreak{}p\allowbreak{}a\allowbreak{}r\allowbreak{}k\allowbreak{}l\allowbreak{}i\allowbreak{}n\allowbreak{}e\allowbreak{}(\allowbreak{})}
  \item
    \texttt{b\allowbreak{}r\allowbreak{}a\allowbreak{}i\allowbreak{}l\allowbreak{}l\allowbreak{}e\allowbreak{}C\allowbreak{}h\allowbreak{}a\allowbreak{}r\allowbreak{}t\allowbreak{}S\allowbreak{}u\allowbreak{}r\allowbreak{}f\allowbreak{}a\allowbreak{}c\allowbreak{}e\allowbreak{}(\allowbreak{})}
  \item
    \texttt{s\allowbreak{}t\allowbreak{}a\allowbreak{}t\allowbreak{}s\allowbreak{}P\allowbreak{}a\allowbreak{}n\allowbreak{}e\allowbreak{}l\allowbreak{}S\allowbreak{}u\allowbreak{}r\allowbreak{}f\allowbreak{}a\allowbreak{}c\allowbreak{}e\allowbreak{}(\allowbreak{})}
  \item
    \texttt{p\allowbreak{}e\allowbreak{}r\allowbreak{}f\allowbreak{}O\allowbreak{}v\allowbreak{}e\allowbreak{}r\allowbreak{}l\allowbreak{}a\allowbreak{}y\allowbreak{}S\allowbreak{}u\allowbreak{}r\allowbreak{}f\allowbreak{}a\allowbreak{}c\allowbreak{}e\allowbreak{}(\allowbreak{})}
  \end{itemize}
\item
  Variants:

  \begin{itemize}
  \tightlist
  \item
    inline trend (sparkline string), high-density area chart (braille Surface), labeled metric panel (boxed key-value Surface), prebuilt performance dashboard (composed Surface)
  \end{itemize}
\item
  Use when:

  \begin{itemize}
  \tightlist
  \item
    the user needs to see numeric trends at a glance without leaving the terminal
  \item
    a time-series or rolling metric is the primary decision input
  \item
    you need a compact performance or status dashboard overlay
  \end{itemize}
\item
  Avoid when:

  \begin{itemize}
  \tightlist
  \item
    exact numeric precision matters more than shape --- use \texttt{t\allowbreak{}a\allowbreak{}b\allowbreak{}l\allowbreak{}e\allowbreak{}(\allowbreak{})} instead
  \item
    the data is categorical rather than continuous --- use \texttt{t\allowbreak{}a\allowbreak{}b\allowbreak{}l\allowbreak{}e\allowbreak{}(\allowbreak{})} or \texttt{t\allowbreak{}r\allowbreak{}e\allowbreak{}e\allowbreak{}(\allowbreak{})}
  \item
    the visualization would be decorative rather than informative
  \end{itemize}
\item
  Ownership:

  \begin{itemize}
  \tightlist
  \item
    Core
  \end{itemize}
\item
  Content guidance:

  \begin{itemize}
  \tightlist
  \item
    use \texttt{s\allowbreak{}p\allowbreak{}a\allowbreak{}r\allowbreak{}k\allowbreak{}l\allowbreak{}i\allowbreak{}n\allowbreak{}e\allowbreak{}(\allowbreak{})} for inline trend hints next to a label; keep it short (8--20 chars)
  \item
    use \texttt{b\allowbreak{}r\allowbreak{}a\allowbreak{}i\allowbreak{}l\allowbreak{}l\allowbreak{}e\allowbreak{}C\allowbreak{}h\allowbreak{}a\allowbreak{}r\allowbreak{}t\allowbreak{}S\allowbreak{}u\allowbreak{}r\allowbreak{}f\allowbreak{}a\allowbreak{}c\allowbreak{}e\allowbreak{}(\allowbreak{})} when the data deserves area-chart density and sub-pixel smoothness
  \item
    use \texttt{s\allowbreak{}t\allowbreak{}a\allowbreak{}t\allowbreak{}s\allowbreak{}P\allowbreak{}a\allowbreak{}n\allowbreak{}e\allowbreak{}l\allowbreak{}S\allowbreak{}u\allowbreak{}r\allowbreak{}f\allowbreak{}a\allowbreak{}c\allowbreak{}e\allowbreak{}(\allowbreak{})} when multiple labeled metrics belong in one titled panel
  \item
    use \texttt{p\allowbreak{}e\allowbreak{}r\allowbreak{}f\allowbreak{}O\allowbreak{}v\allowbreak{}e\allowbreak{}r\allowbreak{}l\allowbreak{}a\allowbreak{}y\allowbreak{}S\allowbreak{}u\allowbreak{}r\allowbreak{}f\allowbreak{}a\allowbreak{}c\allowbreak{}e\allowbreak{}(\allowbreak{})} as the ready-made FPS + memory overlay; prefer it over assembling stats + chart manually
  \item
    always supply \texttt{m\allowbreak{}i\allowbreak{}n}/\texttt{m\allowbreak{}a\allowbreak{}x} when comparing across views so axes stay stable
  \end{itemize}
\item
  Graceful lowering:

  \begin{itemize}
  \tightlist
  \item
    rich: Unicode block or Braille glyphs with semantic color tokens
  \item
    static: same glyph rendering without animation
  \item
    pipe: numeric summary (min/max/mean) or CSV-style output
  \item
    accessible: spoken trend summary (e.g.~``rising from 2 to 8 over 10 samples'')
  \end{itemize}
\item
  Related families:

  \begin{itemize}
  \tightlist
  \item
    \texttt{t\allowbreak{}a\allowbreak{}b\allowbreak{}l\allowbreak{}e\allowbreak{}(\allowbreak{})}
  \item
    \texttt{p\allowbreak{}r\allowbreak{}o\allowbreak{}g\allowbreak{}r\allowbreak{}e\allowbreak{}s\allowbreak{}s\allowbreak{}B\allowbreak{}a\allowbreak{}r\allowbreak{}(\allowbreak{})}
  \item
    \texttt{t\allowbreak{}i\allowbreak{}m\allowbreak{}e\allowbreak{}l\allowbreak{}i\allowbreak{}n\allowbreak{}e\allowbreak{}(\allowbreak{})}
  \end{itemize}
\item
  Carbon analogue:

  \begin{itemize}
  \tightlist
  \item
    sparkline / area chart / metric card
  \end{itemize}
\end{itemize}

\section{Families that still need stronger doctrine}\label{record-019-families-that-still-need-stronger-doctrine}

These are shipped, but the guidance is still thinner than it should be:

\begin{itemize}
\tightlist
\item
  \texttt{s\allowbreak{}e\allowbreak{}p\allowbreak{}a\allowbreak{}r\allowbreak{}a\allowbreak{}t\allowbreak{}o\allowbreak{}r\allowbreak{}(\allowbreak{})} label discipline in dense shells and long documents
\item
  \texttt{t\allowbreak{}a\allowbreak{}b\allowbreak{}s\allowbreak{}(\allowbreak{})} overflow/truncation policy when peer destinations outgrow the available width
\item
  \texttt{w\allowbreak{}i\allowbreak{}z\allowbreak{}a\allowbreak{}r\allowbreak{}d\allowbreak{}(\allowbreak{})} / \texttt{g\allowbreak{}r\allowbreak{}o\allowbreak{}u\allowbreak{}p\allowbreak{}(\allowbreak{})} content architecture for long, interruptible setup flows
\end{itemize}

